\subsection{NILPOTENT LIE GROUPS}

PROPOSITION 12. Let G be a finite-dimensional Lie group. For L(G) to be nilpotent, it is necessary and sufficient that G have a nilpotent open subgroup.

Suppose that \(\mathbf{G}\) has a nilpotent open subgroup \(\mathbf{G}_0\) . By the Corollaries to Propositions 4 and 6, no. 2, \(\mathcal{C}^i\mathrm{L}(G_0) = \{0\}\) for sufficiently large \(i\) . Hence \(\mathrm{L}(G_0) = \mathrm{L}(G)\) is nilpotent.

Suppose that \(\mathrm{L}\langle \mathrm{G}\rangle\) is nilpotent. If \(\mathbf{K} = \mathbf{R}\) or \(\mathbf{C}\) , the identity component \(\mathbf{G}_0\) of \(\mathbf{G}\) is nilpotent by the Corollary to Proposition 4, no. 2, and \(\mathbf{G}_0\) is open in \(\mathbf{G}\) . If \(\mathbf{K}\) is ultrametric, the Corollary to Proposition 6, no. 2, proves that there exist an open subgroup \(\mathbf{G}_1\) of \(\mathbf{G}\) , an integer \(i > 0\) and a neighbourhood \(\mathbf{V}\) of \(e\) in \(\mathbf{G}\) such that \(\mathrm{C}^i\mathrm{G}_1\cap \mathbf{V} = \{\varepsilon \}\) . Then, if \(\mathbf{G}_0\) is a sufficiently small subgroup of \(\mathbf{G}_1\) , \(\mathrm{C}^i\mathrm{G}_0\subset \mathrm{V}\) , hence \(\mathrm{C}^i\mathrm{G}_0 = \{\varepsilon \}\) and \(\mathbf{G}_0\) is nilpotent.

Let \(\mathfrak{g}\) be a nilpotent Lie algebra. The Hausdorff series H(X, Y) corresponding to \(\mathfrak{g}\) has only a finite number of non-zero terms and we know (Chapter II, § 6, no. 5, Remark 3) that the law of composition \((x, y) \mapsto H(x, y)\) defines a group structure on \(\mathfrak{g}\) . Suppose further that \(\mathfrak{g}\) is complete and normable. Clearly the law H is a continuous-polynomial (Differentiable and Analytic Manifolds, R, Appendix). Hence \(\mathfrak{g}\) , together with the law H, is a Lie group G, said to be associated with \(\mathfrak{g}\) . By §4, no. 2, Lemmas 2 and 3, \(\mathrm{L}(\mathrm{G}) = \mathfrak{g}\) . The identity mapping \(\phi\) of \(\mathfrak{g}\) into \(\mathrm{G}\) is an exponential mapping of \(\mathrm{G}\) such that \(\phi(\lambda x)\phi(\lambda' x) = \phi((\lambda + \lambda') x)\) for all \(x \in \mathfrak{g}\) , \(\lambda \in \mathrm{K}\) , \(\lambda' \in \mathrm{K}\) . Every Lie subalgebra \(\mathfrak{h}\) of \(\mathfrak{g}\) admitting a topological supplement is a Lie subgroup \(\mathrm{H}\) of \(\mathrm{G}\) and \(\mathrm{L}(\mathrm{H}) = \mathfrak{h}\) .

PROPOSITION 13. Let G be a finite-dimensional simply connected nilpotent Lie group over R or C.

\begin{enumerate}
\def\labelenumi{(\roman{enumi})}
\item
  \(\exp_{\mathbf{G}}\) is an isomorphism of the Lie group associated with \(\mathrm{L}(\mathrm{G})\) onto \(\mathbf{G}\) .
\item
  Every integral subgroup of \(\mathbf{G}\) is a simply connected Lie subgroup of \(\mathbf{G}\) .
\end{enumerate}

Let \(\mathfrak{g} = \mathrm{L}(\mathrm{G})\) , which is nilpotent (Proposition 12). As two simply connected Lie groups over \(\mathbf{R}\) or \(\mathbf{C}\) which have the same Lie algebra are isomorphic (§ 6, no. 3, Theorem 3 (ii)), it suffices to prove the proposition when \(\mathrm{G}\) is the group associated with \(\mathfrak{g}\) . Then (i) and (ii) follow from what was said before the proposition.

PROPOSITION 14. Let G be a finite-dimensional connected Lie group over R or C.

\begin{enumerate}
\def\labelenumi{(\roman{enumi})}
\item
  If G is nilpotent, \(\exp_{G}\) is étale and surjective.
\item
  If K = C and \(\exp_{G}\) is étale, then G is nilpotent.
\end{enumerate}

Let \(G'\) be the universal covering space of G. Let \(\phi\) be the canonical morphism of \(G'\) onto G. Then \(\exp_G = \phi \circ \exp_{G'}\) (§ 6, no. 4, Proposition 10) and hence (i) follows from Proposition 13 (i).

If \(\mathbf{K} = \mathbf{C}\) and exp is étale, then, for all \(x \in \mathrm{L}(\mathrm{G})\) , \(\operatorname{ad} x\) has no eigenvalue belonging to \(2i\pi (\mathbf{Z} - \{0\})\) (§ 6, no. 4, Corollary to Proposition 12). Applying this to \(\lambda x\) , where \(\lambda\) varies through \(\mathbf{C}\) , it follows that all the eigenvalues of ad \(x\) are zero and hence that ad \(x\) is nilpotent. Therefore \(\mathrm{L}(\mathrm{G})\) is nilpotent (Chapter I, § 4, Corollary 1 to Theorem 1) and hence G is nilpotent (Proposition 12).

PROPOSITION 15. Let \(\mathbf{G}\) be a finite-dimensional connected nilpotent Lie group over \(\mathbf{R}\) or \(\mathbf{C}\) and \(\mathbf{A}\) an integral subgroup of \(\mathbf{G}\) . Then \(Z_{\mathrm{G}}(\mathbf{A})\) is the connected Lie subgroup of \(\mathbf{G}\) with Lie algebra \(\mathfrak{z}_{\mathfrak{g}}(\mathrm{L}(\mathbf{A}))\) .

By Proposition 9 of no. 3 it suffices to prove that \(Z_{\mathrm{G}}(\mathbf{A})\) is connected. Let \(g \in Z_{\mathrm{G}}(\mathbf{A})\) . There exists \(x \in \mathrm{L}(\mathbf{G})\) such that \(g = \exp x\) (Proposition 14). Then \(\operatorname{Ad} g|\mathrm{L}(\mathbf{A}) = 1\) (no. 3, Proposition 9), hence \(\operatorname{Ad} g^n |\mathrm{L}(\mathbf{A}) = 1\) for all \(n \in \mathbf{Z}\) and hence \(\exp (\operatorname{ad} nx)|\mathrm{L}(\mathbf{A}) = 1\) for all \(n \in \mathbf{Z}\) . As the mapping

\[
\lambda \mapsto \exp (\operatorname{ad} \lambda x) | \mathrm{L} (\mathrm{A})
\]

of K into \(\mathcal{L}(\mathrm{L}(\mathrm{A}),\mathrm{L}(\mathrm{G}))\) is polynomial, \(\exp(\operatorname{ad}\lambda x)|\mathrm{L}(\mathrm{A})=1\) for all \(\lambda\in K\) , that is \(\exp(\lambda x)\in Z_{\mathrm{G}}(\mathrm{A})\) for all \(\lambda\in K\) .

PROPOSITION 16. Let G be a finite-dimensional nilpotent Lie group over R or C and A an integral subgroup of G distinct from G. Then \(\mathrm{N}_{\mathrm{G}}(\mathrm{A})\) is a connected Lie subgroup of G distinct from A.

\(N_{\mathrm{G}}(\mathsf{A}) \neq \mathsf{A}\) (Algebra, Chapter I, § 6, Corollary 1 to Proposition 8). By Proposition 11 of no. 4, we need only prove that \(N_{\mathrm{G}}(\mathsf{A})\) is connected. Let \(g \in N_{\mathrm{G}}(\mathsf{A})\) . There exists \(x \in \mathrm{L}(\mathrm{G})\) such that \(g = \exp x\) (Proposition 14). Let \(\mathrm{E}\) be the vector subspace of \(\mathcal{L}(\mathrm{L}(\mathrm{G}))\) consisting of the \(u \in \mathcal{L}(\mathrm{L}(\mathrm{G}))\) such that \(u(\mathrm{L}(\mathsf{A})) \subset \mathrm{L}(\mathsf{A})\) . Then \(\operatorname{Ad} g^n \in \mathbb{E}\) and hence \(\exp (\operatorname{ad} nx) \in \mathbb{E}\) for all \(n \in \mathbf{Z}\) . Hence \(\exp (\operatorname{ad} \lambda x) \in \mathbb{E}\) for all \(\lambda \in K\) , that is \(\exp (\lambda x) \in N_{\mathrm{G}}(\mathsf{A})\) for all \(\lambda \in K\) .

PROPOSITION 17. Let \(\mathfrak{g}\) be a finite-dimensional nilpotent Lie algebra over \(\mathbf{K}\) and \((\mathfrak{g}_0, \mathfrak{g}_1, \ldots, \mathfrak{g}_n)\) a decreasing sequence of ideals of \(\mathfrak{g}\) such that \(\mathfrak{g}_0 = \mathfrak{g}, \mathfrak{g}_n = \{0\}\) and \([\mathfrak{g}, \mathfrak{g}_i] \subset \mathfrak{g}_{i+1}\) for \(0 \leqslant i < n\) . Let \(\mathfrak{a}_1, \mathfrak{a}_2, \ldots, \mathfrak{a}_p\) be vector subspaces of \(\mathfrak{g}\) such that each \(\mathfrak{g}_i\) is the direct sum of its intersections with the \(\mathfrak{a}_j\) . Let \(\mathfrak{g}\) be given the Hausdorff law of composition H. Let \(\phi\) be the mapping

\[
\left(x _ {1}, x _ {2}, \dots , x _ {p}\right) \mapsto x _ {1} \mathsf {H} x _ {2} \mathsf {H} \dots \mathsf {H} x _ {p}
\]

\[
o f a _ {1} \times a _ {2} \times \dots \times a _ {p}
\]

\begin{enumerate}
\def\labelenumi{(\roman{enumi})}
\item
  \(\phi\) is a bijection of \(a_{1} \times a_{2} \times \cdots \times a_{p}\) onto \(g\) ;
\item
  \(\phi\) and \(\phi^{-1}\) are polynomial mappings;
\item
  the mapping \((x,y)\mapsto\phi^{-1}(\phi(x)\cdot\phi(y)^{-1})\) of \((\mathfrak{a}_{1}\times\mathfrak{a}_{2}\times\cdots\times\mathfrak{a}_{p})^{2}\) into \(a_{1}\times a_{2}\times\cdots\times a_{p}\) is polynomial.
\end{enumerate}

The proposition is obvious for \(\dim \mathfrak{g} = 0\) . Suppose that \(\dim \mathfrak{g} > 0\) and that the proposition has been established for dimensions \(<\dim \mathfrak{g}\) . It can be assumed that \(\mathfrak{g}_{n-1} \neq \{0\}\) and \(\mathfrak{g}_{n-1}\) is then a non-zero central ideal of \(\mathfrak{g}\) . There exists an index \(j\) such that \(\mathfrak{b} = \mathfrak{g}_{n-1} \cap \mathfrak{a}_j \neq \{0\}\) . Let \(\mathfrak{g}' = \mathfrak{g}/\mathfrak{b}\) , \(\theta\) be the canonical morphism of \(\mathfrak{g}\) onto \(\mathfrak{g}'\) , \(\mathfrak{g}'_i = \theta(\mathfrak{g}_i)\) and \(\mathfrak{a}_i' = \theta(\mathfrak{a}_i)\) . Then \((\mathfrak{g}_0', \mathfrak{g}_1', \ldots, \mathfrak{g}_n')\) is a decreasing sequence of ideals of \(\mathfrak{g}'\) such that \(\mathfrak{g}_0' = \mathfrak{g}'\) , \(\mathfrak{g}_n' = \{0\}\) ,

\[
[ \mathfrak {g} ^ {\prime}, \mathfrak {g} _ {i} ^ {\prime} ] \subset \mathfrak {g} _ {i + 1} ^ {\prime}
\]

for \(0 \leqslant i < n\) and each \(\mathfrak{g}_i'\) is the direct sum of its intersections with the \(\mathfrak{a}_j'\) . Let \(\phi'\) be the mapping

\[
\left(x _ {1} ^ {\prime}, x _ {2} ^ {\prime}, \dots , x _ {p} ^ {\prime}\right) \mapsto x _ {1} ^ {\prime} \mathsf {H} x _ {2} ^ {\prime} \mathsf {H} \dots \mathsf {H} x _ {p} ^ {\prime}
\]

of \(a_{1}^{\prime} \times a_{2}^{\prime} \times \cdots \times a_{p}^{\prime}\) into \(g^{\prime}\) . By the induction hypothesis, \(\phi^{\prime}\) is bijective and \(\phi^{\prime}, \phi^{\prime-1}\) are polynomial mappings.

Let \(x \in g\) . We write

\[
\phi^ {\prime - 1} (\theta (x)) = \left(x _ {1} ^ {\prime} (x), x _ {2} ^ {\prime} (x), \dots , x _ {p} ^ {\prime} (x)\right).\tag{1}
\]

Then

\[
\theta (x) = x _ {1} ^ {\prime} (x) \mathsf {H} x _ {2} ^ {\prime} (x) \mathsf {H} \dots \mathsf {H} x _ {p} ^ {\prime} (x).\tag{2}
\]

Let \(\mathfrak{h}_1\) be a vector subspace supplementary to \(\mathfrak{b}\) in \(\mathfrak{g}\) , the sum of the \(a_k\) with \(k \neq j\) and a supplement of \(\mathfrak{b}\) in \(a_j\) . There exists a bijection \(\eta\) of \(\mathfrak{g}'\) onto \(\mathfrak{h}_1\) such that \(\theta \circ \eta = \mathrm{Id}_{\mathfrak{g}'}\) . For \(x \in \mathfrak{g}\) , we write

\[
\zeta (x) = \eta (x _ {1} ^ {\prime} (x)) \mathsf {H} \eta (x _ {2} ^ {\prime} (x)) \mathsf {H} \dots \mathsf {H} \eta (x _ {p} ^ {\prime} (x)) \in \mathfrak {g}.\tag{3}
\]

\[
y (x) = \zeta (x) ^ {- 1} \text { H } x = (- \zeta (x))\cdot x.\tag{4}
\]

By (2) and (3), \(\theta (\zeta (x)) = \theta (x)\) and hence \(y(x)\in \mathfrak{h}\) . Finally we write

\[
\psi (x) = \left(\eta \left(x _ {1} ^ {\prime} (x)\right), \dots , \eta \left(x _ {j} ^ {\prime} (x)\right) + y (x), \dots , \eta \left(x _ {p} ^ {\prime} (x)\right)\right) \in a _ {1} \times \dots \times a _ {p}.\tag{5}
\]

As \(y(x)\) is central in \(\mathfrak{g}\) ,

\[
\begin{array}{l l} \phi (\psi (x)) = \eta (x _ {1} ^ {\prime} (x)) \texttt {H} \dots \texttt {H} \eta (x _ {j} ^ {\prime} (x)) \texttt {H} \dots \texttt {H} \eta (x _ {p} ^ {\prime} (x)) \texttt {H} y (x) \\ = \zeta (x) \texttt {H} y (x) & \text {by (3)} \\ = x & \text {by (4)}. \end{array}
\]

Hence \(\phi \circ \psi = \mathrm{Id}_{\mathfrak{g}'}\) . Now let \((x_{1}, x_{2}, \ldots, x_{p}) \in \mathfrak{a}_{1} \times \mathfrak{a}_{2} \times \dots \times \mathfrak{a}_{p}\) and write \(x = \phi(x_{1}, x_{2}, \ldots, x_{p}) = x_{1} \mathsf{H} x_{2} \mathsf{H} \dots \mathsf{H} x_{p}\) . Then

\[
\theta (x) = \theta (x _ {1}) \mathsf {H} \theta (x _ {2}) \mathsf {H} \dots \mathsf {H} \theta (x _ {p}),
\]

hence \(x_{i}^{\prime}(x) = \theta (x_{i})\) for \(1\leqslant i\leqslant p\) and therefore

\[
\begin{array}{c} \zeta (x) = x _ {1} \mathsf {H} x _ {2} \mathsf {H} \dots \mathsf {H} (\eta \theta (x _ {j})) \mathsf {H} \dots \mathsf {H} x _ {p} \\ y (x) = x _ {j} - \eta \theta (x _ {j}). \end{array}
\]

Then by (5)

\[
\psi (x) = \left(x _ {1}, \dots , \eta \theta \left(x _ {j}\right) + x _ {j} - \eta \theta \left(x _ {j}\right), \dots , x _ {p}\right) = \left(x _ {1}, x _ {2}, \dots , x _ {p}\right).
\]

Hence \(\psi \circ \phi = \mathrm{Id}_{\mathfrak{a}_1 \times \ldots \times \mathfrak{a}_p}\) . This proves (i). As the Hausdorff law is polynomial, \(\phi\) is polynomial. By the induction hypothesis, \(\phi'^{-1}\) is polynomial; by formula (1), the functions \(x_j'\) are polynomial, hence \(\zeta\) is polynomial (formula (3)), \(y\) is polynomial (formula (4)) and \(\psi\) is polynomial (formula (5)). This proves (ii). Assertion (iii) follows from (i) and (ii) and the fact that the Hausdorff law is polynomial.

Example of a nilpotent Lie group. Let G be the lower strict triangular subgroup of \(\mathbf{GL}(n,\mathbf{K})\) . It is a Lie subgroup of \(\mathbf{GL}(n,\mathbf{K})\) and \(\mathrm{L}(\mathrm{G})\subset\mathrm{gl}(n,\mathbf{K})\) is the Lie algebra of lower triangular matrices with zero diagonal (§3, no. 10, Proposition 36). By Chapter II, §4, no. 6, Remark, G is nilpotent. Suppose henceforth that K=R or C. As G is homeomorphic to \(\mathbf{K}^{n(n-1)/2}\) , G is simply connected. The exponential mapping of L(G) into G is just the mapping

\[
u \mapsto \exp u = \sum_ {k \geqslant 0} \frac {u ^ {k}}{k !} = \sum_ {k = 0} ^ {n - 1} \frac {u ^ {k}}{k !}
\]

(§ 6, no. 4, Example). By Proposition 13, the exponential is an isomorphism of the manifold L(G) onto the manifold G. Proposition 17 of § 6, no. 9 gives the inverse bijection log. We give \(\mathbf{K}^{n}\) a norm. By Spectral Theories, Chapter I, § 4, no. 9, for g ∈ G and \textbar\textbar g − 1\textbar\textbar{} \textless{} 1,

\[
\log g = \sum_ {k \geqslant 1} \frac {(- 1) ^ {k - 1}}{k} (g - 1) ^ {k}
\]

that is

\[
\log g = \sum_ {k = 1} ^ {n - 1} \frac {(- 1) ^ {k - 1}}{k} (g - 1) ^ {k}.\tag{6}
\]

But the two sides of (6) are analytic functions of \(g\) for \(g \in \mathbf{G}\) and are therefore equal for all \(g \in \mathbf{G}\) .

PROPOSITION 18. Let \(k\) be a commutative field. V a vector space of finite dimension \(>0\) over \(k\) and G a subgroup of \(\mathbf{GL}(\mathbf{V})\) whose elements are unipotent.

\begin{enumerate}
\def\labelenumi{(\roman{enumi})}
\item
  There exists a non-zero element \(v\) of \(V\) such that \(gv = v\) for all \(g \in G\) .
\item
  There exists a basis B of V such that, for all \(g \in G\) , the matrix of \(g\) with respect to B is lower triangular and has all its diagonal elements equal to 1.
\item
  The group G is nilpotent.
\end{enumerate}

\begin{enumerate}
\def\labelenumi{(\alph{enumi})}
\tightlist
\item
  Suppose first that \(k\) is algebraically closed and that the identity representation of \(G\) is simple. Let \(a, b\) be in \(G\) . Then
\end{enumerate}

\[
\operatorname{Tr} (a (b - 1)) = \operatorname{Tr} (a b - 1) - \operatorname{Tr} (a - 1) = 0 - 0 = 0
\]

for \(ab - 1\) and \(a - 1\) are nilpotent. As the vector subspace of \(\mathcal{L}(\mathrm{V})\) generated by \(\mathbf{G}\) is \(\mathcal{L}(\mathrm{V})\) (Algebra, Chapter VIII, §4, Corollary 1 to Proposition 2), \(\operatorname{Tr}(u(b - 1)) = 0\) for all \(u \in \mathcal{L}(\mathrm{V})\) and hence \(b = 1\) . Thus \(\mathrm{G} = \{1\}\) .

\begin{enumerate}
\def\labelenumi{(\alph{enumi})}
\setcounter{enumi}{1}
\item
  We pass now to the general case. Let \(\overline{k}\) be an algebraic closure of \(k\) , \(\overline{\mathrm{V}} = \mathrm{V} \otimes_k \overline{k}\) and \(\overline{\mathbf{G}} \subset \mathbf{GL}(\overline{\mathrm{V}})\) the set of \(a \otimes 1\) for \(a \in \mathbf{G}\) . Let \(\mathrm{W}\) (resp. \(\mathrm{W}'\) ) be the set of elements of \(\mathrm{V}\) (resp. \(\overline{\mathrm{V}}\) ) invariant under \(\mathrm{G}\) (resp. \(\overline{\mathbf{G}}\) ). Then \(\mathrm{W}' = \mathrm{W} \otimes_k \overline{k}\) for \(\mathrm{W} = \bigcap_{g \in \mathbb{G}} \mathrm{Ker}(g - 1)\) and \(\mathrm{W}' = \bigcap_{g \in \mathbb{G}} \mathrm{Ker}(g - 1) \otimes 1\) . If \(\mathrm{V}_1\) denotes a minimal element in the set of non-zero vector subspaces of \(\overline{\mathrm{V}}\) which are stable under \(\overline{\mathbf{G}}\) , then \(\mathrm{V}_1 \subset \mathrm{W}'\) by part (a) of the proof; hence \(\mathrm{W} \neq \{0\}\) , which proves (i).
\item
  By induction on dim V, it follows from (i) that there exists an increasing sequence \(\left(\mathrm{V}_{1},\mathrm{V}_{2},\ldots,\mathrm{V}_{n}\right)\) of vector subspaces of V which are stable under G such that \(V_{n}=V\) and the automorphism group of \(V_{i}/V_{i-1}\) canonically derived from G reduces to \(\{1\}\) for all i (we make the convention that \(V_{r}=\{0\}\) for \(r\leqslant0\) ). This implies first (ii) and consequently (iii) (Chapter II, §4, no. 6, Remark).
\end{enumerate}

COROLLARY 1. Let G be a finite-dimensional connected real or complex Lie group. For G to be nilpotent, it is necessary and sufficient that every element of Ad G be nilpotent.

If every element of Ad G is unipotent, Ad G is nilpotent (Proposition 18) and hence G, which is a central extension of Ad G, is nilpotent. If G is nilpotent, L(G) is nilpotent, hence ad x is nilpotent for all \(x \in \mathrm{L}(G)\) and hence \(\mathrm{Ad}(\exp x) = \exp \mathrm{ad} x\) is unipotent; but every element of G is of the form \(\exp x\) for some \(x \in \mathrm{L}(G)\) (Proposition 14).

COROLLARY 2. Every analytic subgroup of \(\mathbf{GL}(n, \mathbf{K})\) consisting of unipotent elements is a simply connected Lie subgroup.

This follows from Propositions 13 (ii) and 18 (ii) and the fact that the lower strict triangular group is simply connected.

\subsection{SOLVABLE LIE GROUPS}

PROPOSITION 19. Let G be a finite-dimensional Lie group. For L(G) to be solvable, it is necessary and sufficient that G possess a solvable open subgroup.

The proof is analogous to that of Proposition 12 of no. 5.

PROPOSITION 20. Let G be a simply connected solvable Lie group of finite-dimension n over R or C and g = L(G). Let \((\mathfrak{g}_{n}, \mathfrak{g}_{n-1}, \ldots, \mathfrak{g}_{0})\) be a sequence of subalgebras of g of dimensions n, n-1, \ldots, 0, such that \(g_{i-1}\) is an ideal of \(g_{i}\) for i = n, n-1, \ldots, 1.† Let \(G_{i}\) be the integral subgroup of G corresponding to \(g_{i}\) . Let \(x_{i}\) be a vector of \(g_{i}\) not belonging to \(g_{i-1}\) . Let \(\phi_{i}\) be the mapping

\[
\left(\lambda_ {1}, \lambda_ {2}, \dots , \lambda_ {i}\right) \mapsto \left(\exp \lambda_ {1} x _ {1}\right) \left(\exp \lambda_ {2} x _ {2}\right) \dots \left(\exp \lambda_ {i} x _ {i}\right)
\]

of \(\mathbf{K}^i\) into G. Then \(\phi_n\) is an isomorphism of analytic manifolds and \(\phi_i(\mathbf{K}^i) = \mathbf{G}_i\) for all \(i\) .

For \(n = 0\) the proposition is obvious. We argue by induction on \(n\) . Let H be the integral subgroup of G such that \(\mathrm{L(H)} = \mathrm{K}x_n\) . By § 6, no. 6, Corollary 1 to Proposition 14, H and \(G_{n-1}\) are simply connected Lie subgroups of G and, as a Lie group, G is the semi-direct product of H by \(G_{n-1}\) . Hence \(\lambda \mapsto \exp (\lambda x_n)\) is an isomorphism of K onto H and by the induction hypothesis the mapping

\[
\left(\lambda_ {1}, \lambda_ {2}, \dots , \lambda_ {n - 1}\right) \mapsto \left(\exp \lambda_ {1} x _ {1}\right) \left(\exp \lambda_ {2} x _ {2}\right) \dots \left(\exp \lambda_ {n - 1} x _ {n - 1}\right)
\]

is an isomorphism of the analytic manifold \(\mathbf{K}^{n - 1}\) onto the analytic manifold \(\mathbf{G}_{n - 1}\) which maps \(\mathbf{K}^i\times \{0\}\) to \(\mathbf{G}_i\) for \(i = 1,2,\ldots ,n - 1\) . Hence the proposition.

PROPOSITION 21. Let G be a simply connected finite-dimensional solvable Lie group over R or C and M an integral subgroup of G. Then M is a Lie subgroup of G and is simply connected.

We continue to use the notation \(n\) , \(\mathfrak{g}\) , \(\mathfrak{g}_{i}\) , \(x_{i}\) , \(\phi\) of Proposition 20 but impose on the \(x_{i}\) the following supplementary condition: let \(i_{p} > i_{p-1} > \cdots > i_{1}\) be the integers \(i\) such that \(\mathrm{L}(M) \cap \mathfrak{g}_{i} \neq \mathrm{L}(M) \cap \mathfrak{g}_{i-1}\) ; then we take

\[
x _ {i _ {k}} \in \mathrm{L} (\mathrm{M}) \cap \mathfrak {g} _ {i _ {k}}
\]

for \(k = 1, 2, \ldots, p\) . By induction on \(n\) , it is easily seen that \((x_{i_p}, x_{i_{p-1}}, \ldots, x_{i_1})\) is a basis of \(\mathrm{L}(\mathbf{M})\) . Let \(\mathbf{N}\) be a simply connected Lie group such that there exists an isomorphism \(h\) of \(\mathrm{L}(\mathbf{N})\) onto \(\mathrm{L}(\mathbf{M})\) . Let

\[
y _ {p} = h ^ {- 1} (x _ {i _ {p}}), \dots , y _ {1} = h ^ {- 1} (x _ {i _ {1}}).
\]

By Proposition 20, the mapping

\[
\left(\lambda_ {1}, \lambda_ {2}, \dots , \lambda_ {p}\right) \mapsto \left(\exp \lambda_ {1} y _ {1}\right) \left(\exp \lambda_ {2} y _ {2}\right) \dots \left(\exp \lambda_ {p} y _ {p}\right)
\]

is an isomorphism of the manifold \(K^{p}\) onto the manifold N. There exists a Lie group morphism \(\tau\) of N into G such that \(h = L(\tau)\) and \(\tau(N) = M\) ( \(\S 6\) , no. 2, Corollary 1 to Proposition 1). Hence M is the set of elements of G of the form

\[
\begin{array}{r l} \tau ((\exp \lambda_ {1} y _ {1}) \dots (\exp \lambda_ {p} y _ {p})) & = \exp (\lambda_ {1} L (\tau) y _ {1}) \dots \exp (\lambda_ {p} L (\tau) y _ {p}) \\ & = \exp (\lambda_ {1} x _ {i _ {1}}) \dots \exp (\lambda_ {p} x _ {i _ {p}}). \end{array}
\]

Thus \(\mathbf{M} = \phi (\mathbf{T})\) where \(\mathbf{T}\) is a vector subspace of \(\mathbf{K}^n\) .

PROPOSITION 22. Suppose that K = R or C. Let V be a finite-dimensional vector space and G a connected solvable subgroup of GL(V). Suppose that the identity representation of G is simple.

\begin{enumerate}
\def\labelenumi{(\roman{enumi})}
\item
  If \(\mathbf{K} = \mathbf{R}\) , then \(\dim V \leqslant 2\) and \(G\) is commutative.
\item
  If K = C, then dim V = 1.
\item
  Suppose that \(\mathbf{K} = \mathbf{R}\) . Then the closure H of G in \(\mathbf{GL}(V)\) is a solvable connected Lie subgroup of \(\mathbf{GL}(V)\) (no. 1, Corollary 2 to Proposition 1). Hence L(H) is solvable (Proposition 19). The identity representation of L(G) is simple ( \(\S 6\) , no. 5, Corollary 2 to Proposition 13). Hence dim V \(\leqslant 2\) and L(G) is commutative (Chapter I, \(\S 5\) , Corollaries 1 and 4 to Theorem 1). Hence G is commutative.
\item
  Suppose that K = C. Let W be a minimal element among the nonzero real vector subspaces of V which are stable under G. The complex vector subspace of V generated by W is equal to V since the identity representation of G is simple. By (i), G\textbar W is commutative. Hence G is commutative. Therefore every element of G is a homothety (Algebra, Chapter VIII, § 4, Corollary I to Proposition 2), so that dim V = 1.
\end{enumerate}

COROLLARY. Let V be a complex vector space of finite dimension \textgreater0 and G a connected solvable subgroup of GL(V).

\begin{enumerate}
\def\labelenumi{(\roman{enumi})}
\item
  There exists a non-zero element \(v\) of \(V\) such that \(gv \in \mathbf{C}v\) for all \(g \in G\) .
\item
  There exists a basis \(\mathbf{B}\) of \(\mathbf{V}\) such that, for all \(g \in \mathbf{G}\) , the matrix of \(g\) with respect to \(\mathbf{B}\) is lower triangular.
\end{enumerate}

Let \(V_{1}\) be a minimal element among the non-zero vector subspaces of V which are stable under G. By Proposition 22 (ii), \(\dim V_{1} = 1\) . This proves (i). By induction on \(\dim V\) , it then follows that there exists an increasing sequence \((\mathrm{V}_1, \mathrm{V}_2, \ldots, \mathrm{V}_n)\) of vector subspaces of V which are stable under G such that \(\dim \mathrm{V}_{i+1} / \mathrm{V}_i = 1\) for \(i < n\) and \(\mathrm{V}_n = \mathrm{V}\) ; hence (ii).

\subsection{RADICAL OF A LIE GROUP}

PROPOSITION 23. Let G be a finite-dimensional real or complex Lie group, r the radical of L(G) (Chapter I, § 5, Definition 2) and n the largest nilpotent ideal of L(G) (Chapter I, § 4, no. 4). Let R (resp. N) be the integral subgroup of G with Lie algebra r (resp. n). Then R (resp. N) is a solvable (resp. nilpotent) Lie subgroup of G, which is invariant under every continuous automorphism of G. Every solvable (resp. nilpotent) connected normal subgroup of G is contained in R (resp. N).

The group R is solvable (§ 6, Proposition 19). Suppose that K = R. Let G' be a connected solvable normal subgroup of G. Then \(\overline{G'}\) is a connected solvable (no. 1, Corollary 2 to Proposition 1) normal Lie subgroup of G (§ 8, no. 2, Theorem 2). Hence L \((\overline{G'})\) is a solvable ideal of L(G), whence L \((\overline{G'}) \subset r\) and \(\overline{G'} \subset R\) . In particular, \(\overline{R} \subset R\) , hence R is closed and therefore is a Lie subgroup of G. Suppose that K = C. Let H be the underlying real Lie subgroup of G. If \(r'\) is the radical of L(H), \(i'r\) is a solvable ideal of L(H), whence \(r' = i'r'\) ; hence \(r \subset r' \subset r\) and, by the above, R is closed in H and therefore in G; thus R is a Lie subgroup of G. Every connected solvable normal subgroup of G is a connected solvable normal subgroup of H and therefore contained in R. Hence we have proved for K = C and for K = R that R is the largest connected solvable normal subgroup of G; therefore R is invariant under every continuous automorphism of G. The proof for N is completely analogous.

DEFINITION 1. Let G be a finite-dimensional real or complex Lie group. The radical of G is the largest connected solvable normal subgroup of G.

Remark. Even if G is connected there may be solvable normal subgroups of G not contained in the radical of G.

PROPOSITION 24. Suppose that K = R or C. Let \(G_{1}\) , \(G_{2}\) be two finite-dimensional connected Lie groups, \(R_{1}\) and \(R_{2}\) their radicals and \(\phi\) a surjective morphism of \(G_{1}\) into \(G_{2}\) . Then \(\phi(R_{1}) = R_{2}\) .

By § 3, no. 8, Proposition 28, \(\mathrm{L}(\phi)\) is surjective. Hence \(\mathrm{L}(\phi)(\mathrm{L}(\mathbb{R}_1)) = \mathrm{L}(\mathbb{R}_2)\) (Chapter I, § 6, Corollary 3 to Proposition 2). Let \(i\) be the canonical injection of \(\mathbb{R}_1\) into \(\mathbb{G}_1\) . Then the image of \(\phi \circ i\) is \(\mathbb{R}_2\) (§ 6, no. 2, Corollary I to Proposition 1).

PROPOSITION 25. Suppose that \(\mathbf{K} = \mathbf{R}\) or \(\mathbf{C}\) . Let \(\mathbf{G}_1, \mathbf{G}_2\) be finite-dimensional connected Lie groups and \(\mathbf{R}_1\) and \(\mathbf{R}_2\) their radicals. The radical of \(\mathbf{G}_1 \times \mathbf{G}_2\) is \(\mathbf{R}_1 \times \mathbf{R}_2\) . This follows from Chapter I, § 5, Proposition 4.

\subsection{SEMI-SIMPLE LIE GROUPS}

PROPOSITION 26. Let \(G\) be a finite-dimensional connected real or complex Lie group. The following conditions are equivalent:

\begin{enumerate}
\def\labelenumi{(\roman{enumi})}
\item
  L(G) is semi-simple;
\item
  the radical of G is \{e\};
\item
  every normal commutative integral subgroup of \(\mathbf{G}\) is equal to \(\{e\}\) .
\end{enumerate}

Condition (ii) means that the radical of L(G) is \(\{0\}\) and hence (i) \(\Leftrightarrow\) (ii) (Chapter I, § 6, Theorem 1). The equivalence of (i) and (iii) follows from § 6, no. 6, Proposition 14.

DEFINITION 2. A connected real or complex Lie group is called semi-simple if it is finite-dimensional and satisfies the conditions of Proposition 26.

Remark 1. Let G be a finite-dimensional connected real or complex Lie group. If G is not semi-simple, G has a commutative connected Lie subgroup \(G'\) which is invariant under every continuous automorphism, such that \(G' \neq \{\varepsilon\}\) . For let n be the largest nilpotent ideal of L(G); then \(n \neq \{0\}\) and the corresponding analytic subgroup N is a Lie subgroup which is invariant under every continuous automorphism of G (no. 7, Proposition 23); the centre G' of N has the desired properties.

PROPOSITION 27. Let \(G\) be a finite-dimensional connected real or complex Lie group. The following conditions are equivalent:

\begin{enumerate}
\def\labelenumi{(\roman{enumi})}
\item
  L(G) is simple;
\item
  the only normal integral subgroups of \(\mathbf{G}\) are \(\{e\}\) and \(\mathbf{G}\) and further \(\mathbf{G}\) is not commutative.
\end{enumerate}

This follows from § 6, no. 6, Proposition 14.

DEFINITION 3. A connected real or complex Lie group is called almost simple if it is finite-dimensional and satisfies the conditions of Proposition 27.

PROPOSITION 28. Let \(\mathbf{G}\) be a simply connected real or complex Lie group. The following conditions are equivalent:

\begin{enumerate}
\def\labelenumi{(\roman{enumi})}
\item
  G is semi-simple;
\item
  G is isomorphic to the product of a finite number of almost simple groups.
\end{enumerate}

If G is a finite product of almost simple Lie groups, \(\mathrm{L}(\mathrm{G})\) is a finite product of simple Lie algebras and is hence semi-simple. If G is semi-simple, \(\mathrm{L}(\mathrm{G})\) is isomorphic to a product of simple Lie algebras \(L_{1},\ldots,L_{n}\) . Let \(G_{i}\) be a simply connected Lie group with Lie algebra \(L_{i}\) , which is therefore almost simple. Then G and \(G_{1}\times\cdots\times G_{n}\) are simply connected and have isomorphic Lie algebras and are therefore isomorphic.

Lemma 1. Let \(\mathbf{G}\) be a connected topological group, \(\mathbf{Z}\) its centre and \(\mathbf{Z}'\) a discrete subgroup of \(\mathbf{Z}\) . Then the centre of \(\mathbf{G}/\mathbf{Z}'\) is \(\mathbf{Z}/\mathbf{Z}'\) .

Let \(y\) be an element of G whose class modulo \(\mathbf{Z}'\) is a central element of

G/Z'. Let \(\phi\) be the mapping \(g \mapsto gyg^{-1}y^{-1}\) of G into G. Then \(\phi(G)\) is connected and contained in Z' and hence \(\phi(G) = \phi(\{e\}) = \{e\}\) . Therefore \(y \in Z\) .

PROPOSITION 29. Let G be a semi-simple connected real or complex Lie group.

\begin{enumerate}
\def\labelenumi{(\roman{enumi})}
\item
  G = (G, G).
\item
  The centre Z of G is discrete.
\item
  The centre of G/Z is \{e\}.
\end{enumerate}

Assertion (i) follows from the Corollary to Proposition 4, no. 2 and Chapter I, § 6, Theorem 1.

Assertion (ii) follows from § 6, no. 4, Corollary 4 to Proposition 10 and Chapter I, § 6, no. 1, Remark 2.

Assertion (iii) follows from (ii) and Lemma 1.

PROPOSITION 30. (i) Let g be a semi-simple real or complex Lie algebra. Then Int g is the identity component of Aut g.

\begin{enumerate}
\def\labelenumi{(\roman{enumi})}
\setcounter{enumi}{1}
\tightlist
\item
  Let \(\mathbf{G}\) be a semi-simple connected real or complex Lie group. The adjoint group of \(\mathbf{G}\) is the identity component of \(\operatorname{Aut} \mathrm{L}(\mathrm{G})\) . Its centre reduces to the identity element.
\end{enumerate}

Every derivation of g is inner (Chapter I, § 6, Corollary 3 to Proposition 1) and hence L(Int g) = L(Aut g), which proves (i). The first assertion of (ii) follows from (i). The second follows from Proposition 29 (iii) and § 6, no. 4, Corollary 4 (ii) to Proposition 10.

Remark 2. Let g be a complex semi-simple Lie algebra and \(g_{0}\) its underlying real Lie algebra. Then \(\operatorname{Aut}(g)\) is open in \(\operatorname{Aut}(g_{0})\) , for \(\operatorname{Int}(g_{0}) \subset \operatorname{Aut}(g)\) .

PROPOSITION 31. Let G be a simply connected finite-dimensional real or complex Lie group and R its radical. There exists a semi-simple simply connected Lie subgroup S of G such that G, as a Lie group, is the semi-direct product of S by R. If S' is a semi-simple integral subgroup of G, there exists x in the nilpotent radical of L(G) such that

\[
(\operatorname{Ad} \exp x) (\mathrm{S} ^ {\prime}) \subset \mathrm{S}.
\]

This follows from § 6, no. 6, Corollary 1 to Proposition 14 and Chapter I, § 6, Theorem 5 and Corollary 1.

Lemma 2. Let G be a group (resp. a topological group), G' a normal subgroup of G, V a finite-dimensional vector space over a commutative field k (resp. over K), \(\rho\) a linear representation (resp. a continuous linear representation) of G on V and \(\rho' = \rho|G'\) .

\begin{enumerate}
\def\labelenumi{(\roman{enumi})}
\item
  If \(\rho\) is semi-simple, \(\rho'\) is semi-simple.
\item
  If \(\rho'\) is semi-simple and every finite-dimensional linear \(k\) -representation (resp. continuous linear K-representation) of \(\mathrm{G} / \mathrm{G}'\) (resp. \(\mathrm{G} / \overline{\mathrm{G}}'\) ) is semi-simple, then \(\rho\) is semi-simple.
\end{enumerate}

Suppose that \(\rho\) is semi-simple; we prove that \(\rho'\) is semi-simple. It suffices to consider the case where \(\rho\) is simple. Let \(V'\) be a minimal non-zero sub- \(G'\) -module of \(V\) . For all \(g \in G\) , \(\rho(G') \rho(g)V' = \rho(g)\rho(G')V' = \rho(g)V'\) , in other words \(\rho(g)\mathrm{V}'\) is stable under \(\rho(\mathrm{G}')\) ; if \(\mathrm{V}''\) is a sub- \(\mathrm{G}'\) -module of \(\rho(g)\mathrm{V}'\) , then \(\rho(g)^{-1}\mathrm{V}''\) is a sub- \(\mathrm{G}'\) -module of \(\mathrm{V}'\) and hence \(\mathrm{V}''\) is equal to \(\{0\}\) or \(\rho(g)\mathrm{V}'\) . Thus, for all \(g \in \mathrm{G}\) , \(\rho(g)\mathrm{V}'\) is a simple \(\mathrm{G}'\) -module. But \(\sum_{g \in \mathrm{G}} \rho(g)\mathrm{V}'\) is a non-zero sub-G-module of \(\mathrm{V}\) , whence \(\mathrm{V} = \sum_{g \in \mathrm{G}} \rho(g)\mathrm{V}'\) . Hence \(\rho'\) is semi-simple.

Suppose that \(\rho'\) is semi-simple. Let \(W\) be a non-zero sub-G-module of \(V\) . As \(\rho'\) is semi-simple, there exists a projector \(f_0\) of \(V\) onto \(W\) which commutes with \(\rho'(G)\) . Let \(E\) be the set of \(f \in \mathcal{L}(V, V)\) which commute with \(\rho(G')\) , which map \(V\) into \(W\) and whose restriction to \(W\) is a homothety; for \(f \in E\) , let \(\alpha(f)\) denote the ratio of the homothety \(f|W\) . Then \(f_0 \in E\) and \(\alpha(f_0) = 1\) . Clearly \(\alpha\) is a linear form on \(E\) . Let \(F = \text{Ker } \alpha\) , which is a hyperplane of \(E\) . For \(f \in E\) and \(g \in G\) , we write \(\sigma(g)f = \rho(g) \circ f \circ \rho(g)^{-1}\) ; then \(\sigma(g)f\) maps \(V\) into \(W\) and its restriction to \(W\) is the homothety of ratio \(\alpha(f)\) ; if \(g' \in G'\) , then

\[
\begin{array}{r l} \sigma (g) f \circ \rho (g ^ {\prime}) & = \rho (g) \circ f \circ \rho (g) ^ {- 1} \circ \rho (g ^ {\prime}) \\ & = \rho (g) \circ f \circ \rho (g ^ {- 1} g ^ {\prime} g) \circ \rho (g ^ {- 1}) \\ & = \rho (g) \circ \rho (g ^ {- 1} g ^ {\prime} g) \circ f \circ \rho (g ^ {- 1}) \\ & = \rho (g ^ {\prime}) \circ \rho (g) \circ f \circ \rho (g ^ {- 1}) \\ & = \rho (g ^ {\prime}) \circ \sigma (g) f. \end{array}
\]

Hence \(\sigma(g)f\in\mathbf{E}\) . Therefore \(\sigma\) is a linear \(k\) -representation (resp. continuous linear K-representation) of G on E leaving F stable. Then \(\sigma(g)=\mathrm{Id}_{\mathbf{E}}\) for \(g\in\mathbf{G}'\) and hence for \(g\in\overline{\mathbf{G}'}\) in the topological case. Suppose that every finite-dimensional \(k\) -linear representation (resp. continuous K-linear representation) of G/G' (resp. G/ \(\overline{\mathbf{G}'}\) ) is semi-simple. Then there exists in E a supplement of F which is stable under G. In other words, there exists \(f\in\mathbf{E}\) such that \(\alpha(f)=1\) , which is invariant under G. Then \(f\) is a projector of V onto W, and, for \(g\in\mathbf{G}\) , \(\rho(g)\circ f\circ\rho(g^{-1})=f\) , that is \(f\) commutes with \(\rho(\mathbf{G})\) . Thus \(\rho\) is semi-simple.

THEOREM 1. Let G be a finite-dimensional real or complex Lie group, \(G_{0}\) its identity component, R its radical and r the radical of L(G); suppose that \(G/G_{0}\) is finite. Let \(\rho\) be a finite-dimensional analytic linear representation of G. The following conditions are equivalent:

\begin{enumerate}
\def\labelenumi{(\roman{enumi})}
\item
  \(\rho\) is semi-simple;
\item
  \(\rho|G_{0}\) is semi-simple;
\item
  \(\rho|R\) is semi-simple;
\item
  L(ρ) is semi-simple;
\item
  L(ρ)\textbar r is semi-simple.
\end{enumerate}

(i) \(\Leftrightarrow\) (ii) by Lemma 2 and Integration, Chapter VII, §3, Proposition 1. (ii) \(\Leftrightarrow\) (iv) and (iii) \(\Leftrightarrow\) (v) by §6, no. 5, Corollary 2 to Proposition 13. (iv) \(\Leftrightarrow\) (v) by Chapter I, §6, Theorem 4.

COROLLARY 1. Let \(\rho, \rho_{1}, \rho_{2}\) be finite-dimensional semi-simple analytic linear representations of G and n an integer \(\geqslant 0\) . Then \(\rho_{1} \otimes \rho_{2}\) , \(T^{n}\rho\) , \(S^{n}\rho\) , \(\bigwedge^{n}\rho\) (Appendix) are semi-simple.

The semi-simplicity of \(\rho_{1} \otimes \rho_{2}\) follows from Theorem 1 and Chapter I, §6, Corollary 1 to Theorem 4. The semi-simplicity of \(T^{n}_{\rho}, S^{n}_{\rho}, \bigwedge^{n}_{\rho}\) follows from the semi-simplicity of \(\rho_{1} \otimes \rho_{2}\) .

We shall see later that, if \(k\) is a commutative field of characteristic 0, \(\Gamma\) a group and \(\rho_1\) and \(\rho_2\) finite-dimensional semi-simple linear \(k\) -representations of \(\Gamma\) , then \(\rho_1 \otimes \rho_2\) is semi-simple.

COROLLARY 2. Let \(\rho\) be a finite-dimensional semi-simple analytic linear representation of \(G\) on a vector space \(V\) , \(S\) the symmetric algebra of \(V\) and \(S^{\mathfrak{G}}\) the subalgebra of \(S\) consisting of the elements invariant under \((\mathsf{S}(\rho))(G)\) . Then \(S^{\mathfrak{G}}\) is a finitely generated algebra.

This follows from Theorem 1, Chapter I, §6, Theorem 6 (a) and Commutative Algebra, Chapter V, §1, Theorem 2.

COROLLARY 3. Let G be a real or complex Lie group and \(G_{0}\) its identity component. Suppose that \(G_{0}\) is semi-simple and that \(G/G_{0}\) is finite. Then every finite-dimensional analytic linear representation of G is semi-simple.

PROPOSITION 32. Let G be a finite-dimensional connected real Lie group. Suppose that L(G) is reductive. The following conditions are equivalent:

\begin{enumerate}
\def\labelenumi{(\roman{enumi})}
\item
  \(\mathbf{G} / \overline{\mathbf{D}^1}\mathbf{G}\) is compact;
\item
  (resp. (ii')) every finite-dimensional analytic linear representation of \(G\) on a complex (resp. real) vector space is semi-simple.
\item
  \(\Rightarrow\) (ii'): Suppose that \(G / \overline{D}^1 G\) is compact. Then every continuous linear representation of \(G / \overline{D}^1 G\) on a finite-dimensional real vector space is semi-simple (Integration, Chapter VII, §3, Proposition 1). Let \(\rho\) be a finite-dimensional analytic linear representation of \(G\) on a real vector space. Then \(\rho |D^1 G\) is analytic, \(D^1 G\) is semi-simple (Chapter I, §6, Proposition 5) and hence \(\rho |D^1 G\) is semi-simple (Corollary 3 to Theorem 1). Hence \(\rho\) is semi-simple (Lemma 2).
\end{enumerate}

It can be seen similarly that (i) \(\Rightarrow\) (ii).

\begin{enumerate}
\def\labelenumi{(\roman{enumi})}
\tightlist
\item
  \(\Rightarrow\) (ii): Suppose that \(G / \overline{D^1} G\) is not compact and hence isomorphic to a group of the form \(\mathbf{R}^p\times \mathbf{T}^q\) with \(p > 0\) (§6, no. 4, Proposition 11 (ii)). Then there exists a surjective morphism of \(G / \overline{D^1} G\) onto \(\mathbf{R}\) and hence a surjective morphism of \(G\) onto \(\mathbf{R}\) . The mapping
\end{enumerate}

\[
g \mapsto \sigma (g) = \left( \begin{array}{c c} 1 & 0 \\ \rho (g) & 1 \end{array} \right)
\]

is an analytic linear representation of G on \(R^{2}\) which is not semi-simple, for the only 1-dimensional vector subspace of \(\mathbf{R}^2\) which is stable under \(\sigma(G)\) is \(\mathbf{R}(0,1)\) .

It can be seen similarly that (ii) \(\Rightarrow\) (i).

PROPOSITION 33. Let \(G\) be a finite-dimensional complex Lie group whose number of connected components is finite, \(\rho\) a finite-dimensional analytic linear representation of \(G\) and \(G'\) an integral subgroup of the real Lie group \(G\) such that \(L(G')\) generates \(L(G)\) over \(C\) . Then, for \(\rho\) to be semi-simple, it is necessary and sufficient that \(\rho |G'\) be semi-simple.

Let \(\rho' = \rho|G'\) . For \(\rho\) (resp. \(\rho'\) ) to be semi-simple, it is necessary and sufficient that \(L(\rho)\) (resp. \(L(\rho')\) ) be semi-simple (Theorem 1). Let \(V\) be the space of \(\rho\) . For a vector subspace of \(V\) to be stable under \(L(\rho)(L(G))\) , it is necessary and sufficient that it be stable under \(L(\rho')(L(G'))\) . Hence the proposition.

\section{THE AUTOMORPHISM GROUP OF A LIE GROUP}

In this paragraph, K is assumed to be of characteristic zero.

\subsection{INFINITESIMAL AUTOMORPHISMS}

Lemma 1. Let G be a Lie group and \(\alpha\) a vector field on G. For all \(g \in \mathbf{G}\) , let

\[
\beta (g) = \alpha (g) g ^ {- 1} \in \mathrm{L} (\mathrm{G}).
\]

The following conditions are equivalent:

\begin{enumerate}
\def\labelenumi{(\roman{enumi})}
\item
  \(\alpha\) is a homomorphism of the group G into the group T(G);
\item
  for all \(g, g'\) in \(\mathbf{G}\) , \(\alpha(gg') = \alpha(g)g' + g\alpha(g')\) ;
\item
  for all \(g, g'\) in \(\mathbf{G}\) , \(\beta(gg') = \beta(g) + (\mathrm{Ad}g)\beta(g')\) .
\end{enumerate}

Condition (i) means that, for all \(g, g'\) in \(G\) , we have in the group \(T(G)\) :

\[
\beta (g) g \beta (g ^ {\prime}) g ^ {\prime} = \beta (g g ^ {\prime}) g g ^ {\prime}
\]

or

\[
\beta (g) ((\mathrm{Ad} g) \beta (g ^ {\prime})) g g ^ {\prime} = \beta (g g ^ {\prime}) g g ^ {\prime}.
\]

But the product of \(\beta(g)\) and (Ad \(g)\beta(g')\) in T(G) is just the sum of \(\beta(g)\) and (Ad \(g)\beta(g')\) in L(G) (§ 2, no. 1, Proposition 2). Hence (i) \(\Leftrightarrow\) (iii). On the other hand, condition (ii) may be written \(\beta(gg')gg' = \beta(g)gg' + g\beta(g')g'\) , or

\[
\beta (g g ^ {\prime}) = \beta (g) + (\mathrm{Ad} g) \beta (g ^ {\prime})
\]

and hence (ii) \(\Leftrightarrow\) (iii).

DEFINITION 1. Let G be a Lie group. An infinitesimal automorphism of G is any analytic vector field on G satisfying the conditions of Lemma 1.

Lemma 2. Let \(\mathbf{K}'\) be a non-discrete closed subfield of \(\mathbf{K}\) , A a \(\mathbf{K}'\) -manifold, B and C K-manifolds and \(f\) a \(\mathbf{K}'\) -analytic mapping of \(\mathbf{A} \times \mathbf{B}\) into C. Suppose that, for all \(a \in \mathbf{A}\) , the mapping \(b \mapsto f(a, b)\) of \(\mathbf{B}\) into \(\mathbf{C}\) is K-analytic. Then, for all \(t \in \mathbf{TA}\) , the mapping \(u \mapsto (\mathrm{Tf})(t, u)\) of TB into TC is K-analytic.

We fix \(t \in \mathrm{TA}\) and write \(g(u) = (\mathrm{Tf})(t,u)\) . Clearly \(g\) is K'-analytic. By Differentiable and Analytic Manifolds, R, 5.14.6, it suffices to prove that the tangent mappings to \(g\) are K-linear. It can be assumed that A, B, C are open neighbourhoods of 0 in complete normable spaces E, F, G over K', K, K and that \(t\) is tangent to A at 0. We identify TA, TB, TC, with A × E, B × F, C × G and \(t\) with an element of E. Then for all \((x,y) \in \mathrm{TB} = \mathrm{B} \times \mathrm{F}\) .

\[
g (x, y) = \left(f (0, x), \left(\mathrm{D} _ {1} f\right) (0, x) (t) + \left(\mathrm{D} _ {2} f\right) (0, x) (y)\right).
\]

We identify \(\mathrm{T(B\times F)}\) with \((\mathbf{B}\times \mathbf{F})\times (\mathbf{F}\times \mathbf{F})\) and \(\mathrm{T(C\times G)}\) with \((\mathbf{C}\times \mathbf{G})\times (\mathbf{G}\times \mathbf{G})\) . Then, for all

\[
((x, y), (h, k)) \in \mathrm{T} (\mathrm{B} \times \mathrm{F}) = (\mathrm{B} \times \mathrm{F}) \times (\mathrm{F} \times \mathrm{F}),
\]

\((\mathrm{Tg})((x,y),(h,k)) = ((a,b),(c,d)),\) where

\[
\begin{array}{r l} & {a = f (0, x),} \\ & {b = (\mathrm{D} _ {1} f) (0, x) (t) + (\mathrm{D} _ {2} f) (0, x) (y), \qquad c = (\mathrm{D} _ {2} f) (0, x) (h),} \\ & {d = (\mathrm{D} _ {2} \mathrm{D} _ {1} f) (0, x) (t, h) + (\mathrm{D} _ {2} \mathrm{D} _ {2} f) (0, x) (y, h) + (\mathrm{D} _ {2} f) (0, x) (k).} \end{array}
\]

We now fix \((x,y)\in \mathbf{B}\times \mathbf{F}\) . We need to prove that the mapping \((h,k)\mapsto (c,d)\) of \(\mathbf{F}\times \mathbf{F}\) into \(\mathbf{G}\times \mathbf{G}\) is K-linear. As the mapping \(x\mapsto f(0,x)\) of B into C is K-analytic, the mappings

\[
\begin{array}{c} (h, k) \mapsto (\mathrm{D} _ {2} f) (0, x) (h), \qquad (h, k) \mapsto (\mathrm{D} _ {2} \mathrm{D} _ {2} f) (0, x) (y, h), \\ (h, k) \mapsto (\mathrm{D} _ {2} f) (0, x) (k) \end{array}
\]

are K-linear. On the other hand

\[
\left(\mathrm{D} _ {2} \mathrm{D} _ {1} f\right) (0, x) (t, h) = \lim _ {\lambda \in K ^ {*}, \lambda \rightarrow 0} \lambda^ {- 1} \left(\left(\mathrm{D} _ {2} f\right) (\lambda t, x) (h) - \left(\mathrm{D} _ {2} f\right) (0, x) (h)\right)
\]

and, for fixed \(\lambda\) , the mapping \(x \mapsto f(\lambda t, x)\) is K-analytic, so that the mapping \(h \mapsto (\mathrm{D}_2f)(\lambda t, x)(h)\) is K-linear.

PROPOSITION 1. Let \(\mathbf{K}'\) be a non-discrete closed subfield of \(\mathbf{K}\) , \(\mathbf{G}\) a Lie group over \(\mathbf{K}\) , \(\mathbf{V}\) a manifold over \(\mathbf{K}'\) and \((v, g) \mapsto vg\) a \(\mathbf{K}'\) -analytic mapping of \(\mathbf{V} \times \mathbf{G}\) into \(\mathbf{G}\) . Suppose that, for all \(v \in \mathbf{V}\) , the mapping \(g \mapsto vg\) of \(\mathbf{G}\) into \(\mathbf{G}\) is an automorphism of \(\mathbf{G}\) . Let \(\varepsilon\) be an element of \(\mathbf{V}\) such that \(\varepsilon g = g\) for all \(g \in \mathbf{G}\) and \(a \in T_{\varepsilon}(\mathbf{V})\) . Then the vector field \(g \mapsto ag\) on \(\mathbf{G}\) is an infinitesimal automorphism of \(\mathbf{G}\) .

For \(v \in \mathrm{V}\) , \(g_1 \in \mathrm{G}\) , \(g_2 \in \mathrm{G}\) , \(v(g_1g_2) = (vg_1)(vg_2)\) . Hence, for \(u_1 \in \mathrm{TG}\) , \(u_2 \in \mathrm{TG}\) , \(a(u_1u_2) = (au_1)(au_2)\) (§ 2, no. 1, Proposition 3). In particular, the mapping \(g \mapsto ag\) of \(G\) into \(TG\) is a group homomorphism. On the other hand, this mapping is analytic by Lemma 2.

PROPOSITION 2. Let G be a real or complex Lie group and \(\alpha\) an infinitesimal automorphism of G. There exists a law of analytic operation \((\lambda, g) \mapsto \phi_{\lambda}(g)\) of K on G with the following properties:

\begin{enumerate}
\def\labelenumi{(\arabic{enumi})}
\item
  if D is the associated law of infinitesimal operation, then D(1) = α;
\item
  for all \(\lambda \in \mathbf{K}\) , \(\phi_{\lambda} \in \operatorname{Aut} G\) .
\end{enumerate}

\begin{enumerate}
\def\labelenumi{(\alph{enumi})}
\item
  For all \(\mu > 0\) , let \(\mathbf{K}_{\mu}\) be the open ball of centre 0 and radius \(\mu\) in \(\mathbf{K}\) . For all \(g \in G\) , let \(\mathcal{F}_g\) be the set of analytic integral curves \(f\) of \(\alpha\) defined in a ball \(\mathbf{K}_{\mu}\) and such that \(f(0) = g\) . By Differentiable and Analytic Manifolds, R, 9.1.3 and 9.1.5, \(\mathcal{F}_g\) is non-empty and two elements of \(\mathcal{F}_g\) coincide on the intersection of their domains of definition; let \(\mu(g)\) be the least upper bound of the numbers \(\mu\) such that there exists an element of \(\mathcal{F}_g\) defined in \(\mathbf{K}_{\mu}\) ; there exists a unique element of \(\mathcal{F}_g\) defined in \(\mathbf{K}_{\mu(0)}\) ; we denote it by \(f_g\) .
\item
  Let \(g_1, g_2\) be in \(\mathbf{G}\) , \(f_1 \in \mathcal{F}_{g_1}\) , \(f_2 \in \mathcal{F}_{g_2}\) with \(f_1\) and \(f_2\) defined on the same ball \(\mathbf{K}_\mu\) . Then \(f_1f_2: \mathbf{K}_\mu \to \mathbf{G}\) is analytic and \((f_1f_2)(0) = g_1g_2\) . On the other hand, for all \(\lambda \in \mathbf{K}_\mu\) ,
\end{enumerate}

\[
\begin{array}{r l} (\mathrm{T} _ {\lambda} (f _ {1} f _ {2})) 1 & = (\mathrm{T} _ {\lambda} f _ {1}) 1 \cdot f _ {2} (\lambda) + f _ {1} (\lambda) \cdot (\mathrm{T} _ {\lambda} f _ {2}) 1 \quad (\S 2, \text { Proposition   7 }) \\ & = \alpha (f _ {1} (\lambda)) f _ {2} (\lambda) + f _ {1} (\lambda) \alpha (f _ {2} (\lambda)) \\ & = \alpha ((f _ {1} f _ {2}) (\lambda)) \quad \text {(Lemma   1)} \end{array}
\]

and hence \(f_{1}f_{2}\in \mathcal{F}_{g_{1}g_{2}}\) . This proves that \(\mu (g_1g_2)\geqslant \inf (\mu (g_1),\mu (g_2))\)

\begin{enumerate}
\def\labelenumi{(\alph{enumi})}
\setcounter{enumi}{2}
\tightlist
\item
  By Differentiable and Analytic Manifolds, R, 9.1.4 and 9.1.5, there exists a neighbourhood V of \(e\) in G such that \(\sigma = \inf_{g \in \mathbf{V}} \mu(g) > 0\) . Let \(h \in \mathbf{G}\) and C be its connected component. For all \(h' \in \mathbf{C}\) , \(\mu(h') \geqslant \inf(\sigma, \mu(h)) > 0\) by (b). On the other hand, the functions \(f_h'\) , where \(h' \in \mathbf{C}\) , take their values in C. By Differentiable and Analytic Manifolds, R, 9.1.4 and 9.1.5, \(\mu = +\infty\) in C and finally \(\mu = +\infty\) in G. Then let \(f_g(\lambda) = \phi_{\lambda}(g)\) for all \(g \in \mathbf{G}\) and all \(\lambda \in \mathbf{K}\) . By Differentiable and Analytic Manifolds, R, 9.1.4 and 9.1.5, the mapping \((\lambda, g) \mapsto \phi_{\lambda}(g)\) is a law of analytic operation of K on G. Clearly, if D is the associated law of infinitesimal operation, D(1) = \(\alpha\) . By (b),
\end{enumerate}

\[
\phi_ {\lambda} (g _ {1} g _ {2}) = \phi_ {\lambda} (g _ {1}) \phi_ {\lambda} (g _ {2})
\]

for all \(\lambda \in \mathbf{K}\) , \(g_{1} \in \mathbf{G}\) , \(g_{2} \in \mathbf{G}\) .

PROPOSITION 3. Suppose that K is ultrametric. Let G be a compact Lie group and \(\alpha\) an infinitesimal automorphism of G. There exist an open subgroup I of K and a law of analytic operation \((\lambda, g) \mapsto \phi_{\lambda}(g)\) of I on G with the following properties:

\begin{enumerate}
\def\labelenumi{(\arabic{enumi})}
\item
  if D is the associated law of infinitesimal operation, then D(1) = α;
\item
  for all \(\lambda \in \mathbf{I}\) , \(\phi_{\lambda} \in \operatorname{Aut} G\) .
\end{enumerate}

As G is compact, there exist an open subgroup I' of K and a law of analytic operation \((\lambda, g) \mapsto \phi_{\lambda}(g)\) of I' on G with property (1) of the proposition (§ 4, no. 7, Corollary 2 to Theorem 6). We write \(\phi_{\lambda}(g) = f_{g}(\lambda)\) for \(\lambda \in I'\) and \(g \in G\) . Then, for \(g_{1}, g_{2}\) in \(G\) and \(\lambda \in I'\) ,

\[
\begin{array}{r l} (\mathrm{T} _ {\lambda} (f _ {g _ {1}} f _ {g _ {2}})) 1 & = (\mathrm{T} _ {\lambda} f _ {g _ {1}}) 1 \cdot f _ {g _ {2}} (\lambda) + f _ {g _ {1}} (\lambda) \cdot (\mathrm{T} _ {\lambda} f _ {g _ {2}}) 1 \\ & = \alpha (f _ {g _ {1}} (\lambda)) f _ {g _ {2}} (\lambda) + f _ {g _ {1}} (\lambda) \alpha (f _ {g _ {2}} (\lambda)) \\ & = \alpha (f _ {g _ {1}} (\lambda) f _ {g _ {2}} (\lambda)) \end{array}
\]

and \((f_{g_1}f_{g_2})(0) = g_1g_2 = f_{g_1g_2}(0)\) . Hence \(f_{g_1'g_2'}(\lambda) = f_{g_1'}(\lambda)f_{g_2'}(\lambda)\) for

\[
(g _ {1} ^ {\prime}, g _ {2} ^ {\prime}, \lambda)
\]

in a neighbourhood of \((g_1, g_2, 0)\) (Differentiable and Analytic Manifolds, R, 9.1.8). As G is compact, there exists an open subgroup I of I' such that \(f_{g_1 g_2}(\lambda) = f_{g_1}(\lambda) f_{g_2}(\lambda)\) for all \(g_1 \in G\) , \(g_2 \in G\) , \(\lambda \in I\) . In other words, \(\phi_\lambda \in \operatorname{Aut} G\) for \(\lambda \in I\) .

Lemma 3. Let G and G' be Lie groups and \(\phi\) a homomorphism of G into \(\operatorname{Aut}(\mathbf{G}')\) . Let \(f(g, g') = (\phi(g))(g')\) for \(g \in \mathbf{G}\) , \(g' \in \mathbf{G}'\) . Consider the following conditions:

\begin{enumerate}
\def\labelenumi{(\roman{enumi})}
\item
  \(f\) is analytic;
\item
  \(f\) is analytic in a neighbourhood of \((e_{\mathbf{G}}, e_{\mathbf{G}^{\prime}})\) ;
\item
  for all \(g' \in G'\) , the mapping \(g \mapsto f(g, g')\) is analytic.
\end{enumerate}

Then (i) \(\Leftrightarrow\) ((ii) and (iii)). If G is connected, (i) \(\Leftrightarrow\) (ii).

Clearly (i) implies (ii) and (iii). Let \(g_0 \in G, g_0' \in G'\) . For all \(g \in G, g' \in G'\) ,

\[
f \left(g g _ {0}, g ^ {\prime} g _ {0} ^ {\prime}\right) = (\phi (g) \phi (g _ {0})) \left(g ^ {\prime} g _ {0} ^ {\prime}\right) = \phi (g) \left(\phi (g _ {0}) g ^ {\prime}\right)\cdot \phi (g) \left(\phi (g _ {0}) g _ {0} ^ {\prime}\right).
\]

This proves the implication ((ii) and (iii)) \(\Rightarrow\) (i). Finally, if \(\mathbf{G}'\) is connected, \(\mathbf{G}'\) is generated by every neighbourhood of \(e_{\mathbf{G}'}\) and hence (ii) \(\Rightarrow\) (iii).

\subsection{THE AUTOMORPHISM GROUP OF A LIE GROUP (REAL OR COMPLEX CASE)}

In this no., we assume that K = R or C.

Lemma 4. Let \(\mathbf{H}\) be a finite-dimensional simply connected Lie group.

\begin{enumerate}
\def\labelenumi{(\roman{enumi})}
\item
  For all \(u \in \operatorname{Aut} \mathbf{L}(\mathbf{H})\) , let \(\theta(u)\) be the unique automorphism of \(\mathbf{H}\) such that \(\mathbf{L}(\theta(u)) = u\) . Then the mapping \((u, g) \mapsto \theta(u)g\) of \((\operatorname{Aut} \mathbf{L}(\mathbf{H})) \times \mathbf{H}\) into \(\mathbf{H}\) is analytic.
\item
  Let \(\mathbf{N}\) be a Lie subgroup of \(\mathbf{H}\) and \(\operatorname{Aut}(\mathbf{H},\mathbf{N})\) the set of \(v\in \operatorname{Aut}\mathbf{H}\) such that \(v(\mathbf{N}) = \mathbf{N}\) . Then \(\theta^{-1}(\operatorname {Aut}(\mathbf{H},\mathbf{N}))\) is a Lie subgroup of \(\operatorname {Aut}\mathbf{L}(\mathbf{H})\) .
\item
  Suppose that N is discrete and normal, so that the Lie algebra of G = H/N is identified with L(H). For all \(w \in Aut G\) , let \(\eta(w)\) be the unique automorphism of H such that \(\mathrm{L}(\eta(w)) = \mathrm{L}(w)\) . Then the mapping \(\eta\) is an isomorphism of the group Aut G onto the group \(\mathrm{Aut}(\mathrm{H}, \mathrm{N})\) .
\end{enumerate}

To prove (i), it suffices, by Lemma 3 of no. 1, to verify that the mapping \((u,g)\mapsto\theta(u)g\) is analytic in a neighbourhood of \((\mathrm{Id}_{\mathrm{L(H)}},e)\) . There exists an open neighbourhood B of 0 in L(H) such that \(\psi = \exp_{\mathbf{H}}|\mathbf{B}\) is an analytic isomorphism of B onto an open neighbourhood of \(e\) in H. There exist an open neighbourhood U of \(\mathrm{Id}_{\mathrm{L(H)}}\) in Aut L(H) and an open neighbourhood B' of 0 in L(H) such that \(\mathrm{U(B')} \subset \mathrm{B}\) . Then the mapping \((u, g) \mapsto \theta(u)g\) of \(\mathrm{U} \times \psi(\mathrm{B'})\) into H is composed of the following mappings:

\[
\begin{array}{l l l l} \text {the mapping} (u, g) \mapsto \langle u, \psi^ {- 1} (g) \rangle & \text {of} & \mathrm{U} \times \psi (\mathrm{B} ^ {\prime}) & \text {into} \quad \mathrm{U} \times \mathrm{B} ^ {\prime}; \\ \text {the mapping} (u, x) \mapsto u (x) & \text {of} & \mathrm{U} \times \mathrm{B} ^ {\prime} & \text {into} \quad \mathrm{B}; \\ \text {the mapping} y \mapsto \psi (y) & \text {of} & \mathrm{B} & \text {into} \quad \mathrm{G}. \end{array}
\]

Hence this mapping is analytic.

Let \(p\) be the canonical mapping of \(H\) into the homogeneous space \(H/N\) . Then \(\theta^{-1}(\text{Aut}(H, N))\) is the set of \(u \in \text{Aut } L(H)\) such that

\[
p (\theta (u) g) = p (e), \quad p (\theta (u ^ {- 1}) g) = p (e)
\]

for all \(g \in \mathbb{N}\) . By § 8, no. 2, Theorem 2 and Corollary 2 to Theorem 2, this proves (ii).

Suppose that N is discrete and normal. Let \(w \in \operatorname{Aut} G\) . Then

\[
\mathrm{L} (p \circ \eta (w)) = \mathrm{L} (\eta (w)) = \mathrm{L} (w) = \mathrm{L} (w \circ p)
\]

hence \(p \circ \eta(w) = w \circ p\) and therefore \(\eta(w) \in \text{Aut(H, N)}\) . Clearly the mapping \(\eta\) of Aut G into Aut(H, N) is an injective homomorphism. This homomorphism is surjective because \(p: H \to G\) is a submersion.

Let G be a locally compact group and \(\Gamma\) the automorphism group of G. Recall that a topology \(T_{\beta}\) has been defined on \(\Gamma\) (General Topology, Chapter X, §3, no. 5). It is the coarsest topology for which the mappings \(v \mapsto v\) and \(v \mapsto v^{-1}\) of \(\Gamma\) into \(\mathcal{C}_{\varepsilon}(G; G)\) (space of continuous mappings of G into G with the compact convergence topology) are continuous. The topology \(T_{\beta}\) is compatible with the group structure on \(\Gamma\) (loc. cit.). For every compact subset L of G and every neighbourhood U of \(e_{\mathrm{G}}\) in G, let N(L, U) be the set of \(\phi \in \Gamma\) such that \(\phi(g) \in gU\) and \(\phi^{-1}(g) \in gU\) for all \(g \in L\) ; then the N(L, U) form a fundamental system of neighbourhoods of \(e_{\Gamma}\) . If G is generated by a compact subset C, the topology \(T_{\beta}\) is also the coarsest topology for which the mappings \(v \mapsto v[C\) and \(v \mapsto v^{-1}[C]\) of \(\Gamma\) into \(\mathcal{C}_{u}(C; G)\) are continuous (for every compact subset of G is contained in \((\mathrm{C} \cup \mathrm{C}^{-1})^{\mathrm{n}}\) for sufficiently large n). If K is locally compact and V is a finite-dimensional vector space over K, the topology \(T_{\beta}\) on GL(V) is just the usual topology.

THEOREM 1. Let G be a finite-dimensional Lie group and \(G_{0}\) its identity component. Suppose that G is generated by \(G_{0}\) and a finite number of elements.

\begin{enumerate}
\def\labelenumi{(\roman{enumi})}
\tightlist
\item
  There exists on Aut G one and only one analytic manifold structure satisfying the following condition:
\end{enumerate}

(AUT) for every analytic manifold M and every mapping f of M into Aut G, f is analytic if and only if the mapping \((m, g) \mapsto f(m)g\) of M × G into G is analytic.

Suppose in the rest of the statement that Aut G has this structure.

\begin{enumerate}
\def\labelenumi{(\roman{enumi})}
\setcounter{enumi}{1}
\item
  Aut G is a finite-dimensional Lie group.
\item
  The morphism \(\phi: u \mapsto L(u)\) of Aut G into Aut L(G) is analytic.
\item
  If G is connected, \(\phi\) is an isomorphism of the Lie group Aut G onto a Lie subgroup of Aut L(G); this Lie subgroup is equal to Aut L(G) if G is simply connected.
\item
  Let \(\mathfrak{a}\) be the set of infinitesimal automorphisms of \(\mathbf{G}\) . Then \(\mathfrak{a}\) is a Lie algebra of vector fields and the law of infinitesimal operation associated with the mapping \((u, g) \mapsto u(g)\) of (Aut G) \(\times\) G into G is an isomorphism of L(Aut G) onto \(\mathfrak{a}\) .
\item
  The topology of the Lie group Aut G is the topology \(\mathcal{T}_{\beta}\) .
\end{enumerate}

\begin{enumerate}
\def\labelenumi{(\alph{enumi})}
\item
  The uniqueness of the analytic structure considered in (i) is obvious.
\item
  Suppose that G is connected. Let H be the universal covering space of G, p the canonical morphism of H onto G and N = Ker p.~We introduce the notation \(\theta\) , \(\eta\) and \(\operatorname{Aut}(\mathrm{H},\mathrm{N})\) of Lemma 4. We transport the Lie group structure of Aut L(G) to Aut H by means of \(\theta\) . Then Aut H becomes a finite-dimensional Lie group and \(\operatorname{Aut}(\mathrm{H},\mathrm{N})\) a Lie subgroup of Aut H (Lemma 4 (ii)). We transport the Lie group structure of \(\operatorname{Aut}(\mathrm{H},\mathrm{N})\) to Aut G by means of \(\eta^{-1}\) . Then Aut G becomes a finite-dimensional Lie group. Properties (ii), (iii) and (iv) of the theorem are satisfied and the mapping \((u,g)\mapsto u(g)\) of \((\operatorname{Aut}\mathbf{G})\times\mathbf{G}\) into G is analytic (Lemma 4 (i)). Let M be an analytic manifold, f a mapping of M into Aut G and \(\phi\) the mapping \((m,g)\mapsto f(m)g\) of M × G into G. Clearly, if f is analytic, \(\phi\) is analytic. Suppose that \(\phi\) is analytic. Then the T \(\phi\) : TM × TG → TG is analytic; its restriction to M × L(G), that is the mapping \((m,x)\mapsto\operatorname{L}(f(m))x\) of M × L(G) into L(G) is therefore analytic; as L(G) is finite-dimensional, it follows that the mapping \(m\mapsto\operatorname{L}(f(m))\) of M into Aut L(G) is analytic and hence that f is analytic. Thus (i) holds.
\end{enumerate}

Let \(\mathrm{L}(\mathrm{G})\) be given a norm. For all \(\lambda > 0\) , let \(\mathrm{B}_{\lambda}\) be the open ball of centre 0 and radius \(\lambda\) in \(\mathrm{L}(\mathrm{G})\) . We choose \(\lambda > 0\) sufficiently small for \(\psi = \exp_{\mathrm{G}}|\mathrm{B}_{\lambda}\) to be an isomorphism of the analytic manifold \(\mathrm{B}_{\lambda}\) onto the open submanifold \(\psi(\mathrm{B}_{\lambda})\) of \(G\) . Let \(\Phi\) be a filter on Aut \(G\) . For \(\Phi\) to converge to \(\mathrm{Id}_{\mathrm{G}}\) in Aut \(G\) , it is necessary and sufficient that \(\mathrm{L}(\Phi)\) converge to \(\mathrm{Id}_{\mathrm{L}(\mathrm{G})}\) in Aut \(L(\mathrm{G})\) and hence that \(L(\Phi)|B_{\lambda/2}\) and \(L(\Phi)^{-1}\mathrm{B}_{\lambda/2}\) converge uniformly to \(\mathrm{Id}_{\mathrm{B}_{\lambda/2}}\) . This condition implies that \(\Phi|\psi(\mathrm{B}_{\lambda/2})\) and \(\Phi^{-1}|\psi(\mathrm{B}_{\lambda/2})\) converge uniformly to \(\mathrm{Id}_{\psi(\mathrm{B}_{\lambda/2})}\) . Conversely, suppose that \(\Phi|\psi(\mathrm{B}_{\lambda/2})\) converges uniformly to \(\mathrm{Id}_{\psi(\mathrm{B}_{\lambda/2})}\) . There exists \(M \in \Phi\) such that, if \(u \in M\) , then \(u(\psi(\mathrm{B}_{\lambda/2})) \subset \psi(\mathrm{B}_{2\lambda/3})\) ; then \(L(u)(\mathrm{B}_{\lambda/2})\) is a connected subset of \(L(\mathrm{G})\) whose image under \(\exp_{\mathrm{G}}\) is contained in \(\psi(\mathrm{B}_{2\lambda/3})\) , hence \(L(u)(\mathrm{B}_{\lambda/2})\) does not meet \(B_{\lambda} - B_{2\lambda/3}\) and therefore \(L(u)(\mathrm{B}_{\lambda/2}) \subset B_{\lambda}\) ; then the hypothesis that \(\Phi|\psi(\mathrm{B}_{\lambda/2})\) converges uniformly to \(\mathrm{Id}_{\psi(\mathrm{B}_{\lambda/2})}\) implies that \(L(\Phi)|B_{\lambda/2}\) converges uniformly to \(\mathrm{Id}_{\mathrm{B}_{\lambda/2}}\) . It then follows that:

\((\Phi\) converges to \(\mathrm{Id}_{\mathrm{G}}\) in Aut G) \(\Leftrightarrow\) ( \(\Phi\) converges to \(\mathrm{Id}_{\mathrm{G}}\) under \(\mathcal{T}_{\beta}\) ).

This proves (vi).

Let D be the law of infinitesimal operation associated with the law of left operation of \(\mathrm{Aut(G)}\) on G. By Propositions 1 and 2 of no. 1, \(\mathrm{D(L(AutG))} = a\) . Hence \(a\) is a Lie algebra of vector fields and D is a morphism of \(\mathrm{L(AutG)}\) onto \(a\) . Let \(x_{1}\) and \(x_{2}\) be elements of \(\mathrm{L(AutG)}\) such that \(\mathrm{D}(x_1) = \mathrm{D}(x_2)\) . Then the laws of operation \((\lambda, g) \mapsto (\exp \lambda x_1)g\) and \((\lambda, g) \mapsto (\exp \lambda x_2)g\) of K on G have the same associated law of infinitesimal operation; hence, for \(|\lambda|\) sufficiently small, \(\exp \lambda x_1\) and \(\exp \lambda x_2\) coincide on a neighbourhood of \(e\) (§ 4, no. 7, Theorem 6), whence \(\exp \lambda x_1 = \exp \lambda x_2\) . It follows that \(x_1 = x_2\) and hence D is an isomorphism of \(\mathrm{L(AutG)}\) onto \(a\) .

The theorem has thus been completely proved for G connected.

\begin{enumerate}
\def\labelenumi{(\alph{enumi})}
\setcounter{enumi}{2}
\tightlist
\item
  We pass to the general case. By hypothesis, G is generated by \(G_{0}\) and a finite number of elements \(x_{1}, x_{2}, \ldots, x_{n}\) . Every \(u \in AutG\) leaves \(G_{0}\) stable. Let \(Aut_{1}G\) be the set of \(u \in AutG\) which, on passing to the quotient, give the identity automorphism of \(G/G_{0}\) . This is a normal subgroup of Aut G. By part (b) of the proof, \(AutG_{0}\) has a canonical Lie group structure and the mapping \((g_{1}, g_{2}, \ldots, g_{n}, u) \mapsto (ug_{1}, ug_{2}, \ldots, ug_{n})\) of \(G_{0}^{n} \times AutG_{0}\) into \(G_{0}^{n}\) is analytic. Let P be the corresponding semidirect product of \(AutG_{0}\) by \(G_{0}^{n}\) ; it is a finite-dimensional Lie group ( \(\S\) 1, no. 4, Proposition 7).
\end{enumerate}

If \(w\in \mathrm{Aut}_1\mathrm{G}\) , we write

\[
\begin{array}{r l} w _ {0} & = w | \mathrm{G} _ {0} \in \mathrm{Aut}   \mathrm{G} _ {0} \\ w _ {i} & = x _ {i} ^ {- 1} w (x _ {i}) \in \mathrm{G} _ {0} \quad (1 \leqslant i \leqslant n) \\ \zeta (w) & = ((w _ {1}, \ldots , w _ {n}), w _ {0}) \in \mathrm{P}. \end{array}
\]

For all \(w, w'\) in \(\mathrm{Aut}_1\mathrm{G}\) ,

\[
\begin{array}{r l} \zeta (w) \zeta (w ^ {\prime}) & = \left((w _ {1}, \ldots , w _ {n}), w _ {0}\right) \left((w _ {1} ^ {\prime}, \ldots , w _ {n} ^ {\prime}), w _ {0} ^ {\prime}\right) \\ & = \left((w _ {1} w _ {0} (w _ {1} ^ {\prime}), \ldots , w _ {n} w _ {0} (w _ {n} ^ {\prime})), w _ {0} w _ {0} ^ {\prime}\right) \\ & = \left((x _ {1} ^ {- 1} w (x _ {1}) w (x _ {1} ^ {- 1} w ^ {\prime} (x _ {1})), \ldots , x _ {n} ^ {- 1} w (x _ {n}) w (x _ {n} ^ {- 1} w ^ {\prime} (x _ {n}))), w _ {0} w _ {0} ^ {\prime}\right) \\ & = \left(((w w ^ {\prime}) _ {1}, \ldots , (w w ^ {\prime}) _ {n}), (w w ^ {\prime}) _ {0}\right) \\ & = \zeta (w w ^ {\prime}) \end{array}
\]

and hence \(\zeta\) is a homomorphism of \(\mathrm{Aut}_1\mathrm{G}\) into P. This homomorphism is obviously injective.

We show that \(\zeta(\mathrm{Aut}_1\mathrm{G})\) is closed in P. Let \(\Phi\) be a filter on \(\mathrm{Aut}_1\mathrm{G}\) such that \(\zeta(\Phi)\) converges to a point \(((w_1,\dots,w_n),w_0)\) of P. Then \(\Phi\) converges pointwise to a mapping \(v\) of G into G. Clearly \(v\) is an endomorphism of the group G. Moreover, \(v\) leaves each coset modulo \(\mathbf{G}_0\) stable and \(v |\mathrm{G}_0 = w_0\) . It follows that \(v\in \mathrm{Aut}_1\mathrm{G}\) . As \(\zeta(v) = ((w_1,\dots,w_n),w_0)\) , we have shown that \(\zeta(\mathrm{Aut}_1\mathrm{G})\) is closed in P.

\begin{enumerate}
\def\labelenumi{(\alph{enumi})}
\setcounter{enumi}{3}
\tightlist
\item
  In part (d) of the proof we assume that \(\mathbf{K} = \mathbf{R}\) . By § 8, no. 2, Theorem 2, \(\zeta(\mathrm{Aut}_1\mathrm{G})\) is a Lie subgroup of P. We transport the real Lie group structure on \(\zeta(\mathrm{Aut}_1\mathrm{G})\) to \(\mathrm{Aut}_1\mathrm{G}\) by means of \(\zeta^{-1}\) . Thus \(\mathrm{Aut}_1\mathrm{G}\) becomes a finite-dimensional Lie group.
\end{enumerate}

Let M be an analytic manifold, f a mapping of M into \(Aut_{1}G\) and \(\phi\) the mapping \((m,g)\mapsto f(m)g\) of \(M\times G\) into G. We have the following equivalences:

\(f\) analytic \(\Leftrightarrow\) the mappings \(m\mapsto (f(m))_i\) , where \(0\leqslant i\leqslant n\) , are analytic \(\Leftrightarrow \left\{ \begin{array}{ll}\text{the mappings } m\mapsto f(m)x_i\text{ of M into G, for } 1\leqslant i\leqslant n,\text{ are analytic} \\ \text{and} \\ \text{the mapping } (m,g)\mapsto f(m)g\text{ of M} \times \mathrm{G}_0\text{ into G is analytic} \end{array} \right.\) \(\Leftrightarrow \phi\) is analytic.

For \(w \in \mathrm{Aut}_1\mathrm{G}\) , \(\mathrm{L}(w) = \mathrm{L}(w_0)\) and hence the morphism \(w \mapsto \mathrm{L}(w)\) of \(\mathrm{Aut}_1\mathrm{G}\) into \(\mathrm{Aut}\, \mathrm{L}(\mathrm{G})\) is analytic. We see as in (b) that the law of infinitesimal operation associated with the law of operation of \(\mathrm{Aut}_1\mathrm{G}\) on \(\mathrm{G}\) is an isomorphism of \(\mathrm{L}(\mathrm{Aut}_1\mathrm{G})\) onto \(\mathfrak{a}\) .

Let C be a compact subset of \(G_{0}\) generating \(G_{0}\) . For a filter \(\Phi\) to converge to \(Id_{G}\) on \(Aut_{1}G\) , it is necessary and sufficient that \(\Phi|(C \cup \{x_{1}\} \cup \cdots \cup \{x_{n}\})\) and \(\Phi^{-1}|(C \cup \{x_{1}\} \cup \cdots \cup \{x_{n}\})\) converge uniformly to

\[
\operatorname{Id} _ {\mathbb {G}} \left| \left(\mathrm{C} \cup \{x _ {1} \} \cup \dots \cup \{x _ {n} \}\right). \right.
\]

The topology of \(\mathrm{Aut}_1\mathrm{G}\) is therefore the topology \(\mathcal{T}_{\beta}\) .

Clearly \(\mathrm{Aut}_1\mathrm{G}\) is open in Aut G with the topology \(\mathcal{T}_{\beta}\) . There exists on Aut G a Lie group structure compatible with this topology and inducing on \(\mathrm{Aut}_1\mathrm{G}\) the structure constructed above (§ 8, no. 1, Corollary 2 to Theorem 1). The fact that the Lie group Aut G has the properties of the theorem follows from the corresponding properties for \(\mathrm{Aut}_1\mathrm{G}\) .

\begin{enumerate}
\def\labelenumi{(\alph{enumi})}
\setcounter{enumi}{4}
\tightlist
\item
  In part (e) of the proof we assume that \(\mathbf{K} = \mathbf{C}\) . By (c) and Theorem 2 of §8, no. 2, there exists on \(\mathrm{Aut}_1\mathrm{G}\) a real Lie group structure such that \(\zeta\) is an isomorphism of \(\mathrm{Aut}_1\mathrm{G}\) onto a real Lie subgroup of \(\mathbb{P}\) .
\end{enumerate}

The law of operation \((w, g) \mapsto wg\) of \((\mathrm{Aut}_1\mathrm{G}) \times \mathrm{G}\) on \(\mathbf{G}\) is real analytic. Let \(\mathbf{D}\) be the associated law of infinitesimal operation. By Propositions 1 and 2 of no. 1, \(\mathrm{D}(\mathrm{L}(\mathrm{Aut}_1\mathrm{G})) = \mathfrak{a}\) .

For all \(\alpha \in \mathfrak{a}\) , let \(\alpha_0\) denote the restriction of \(\alpha\) to \(\mathrm{G}_0\) ; it is an infinitesimal automorphism of \(\mathrm{G}_0\) which we identify, because of part (b) of the proof, with an element of \(\mathrm{L}(\mathrm{Aut}\mathrm{G}_0)\) . For \(1 \leqslant i \leqslant n\) , we write

\[
\alpha_ {i} = x _ {i} ^ {- 1} \alpha (x _ {i}) \in \mathrm{L} (\mathrm{G}) = \mathrm{L} (\mathrm{G} _ {0}).
\]

Finally, we write \(f(\alpha) = ((\alpha_1, \ldots, \alpha_n), \alpha_0) \in \mathrm{L}(\mathrm{P})\) . Then \(f\) is a \(\mathbf{C}\) -linear mapping of \(\mathfrak{a}\) into \(\mathrm{L}(\mathrm{P})\) .

On the other hand, clearly \(\mathrm{L}(\zeta) = f\circ \mathrm{D}\) . Hence \(\mathrm{L}(\zeta)(\mathrm{L}(\mathrm{Aut}_1\mathrm{G})) = f(a)\) is a complex vector subspace of \(\mathrm{L}(P)\) . By Proposition 2 of § 4, no. 2, \(\zeta (\mathrm{Aut}_1\mathrm{G})\) is a complex Lie subgroup of \(P\) and we can proceed exactly as in (d): we transport the complex Lie group structure on \(\zeta (\mathrm{Aut}_1\mathrm{G})\) to \(\mathrm{Aut}_1\mathrm{G}\) by means of \(\zeta^{-1}\) and we see as in (d) that \(\mathrm{Aut}_1\mathrm{G}\) has the properties analogous to properties (i), (ii), (iii), (v) and (vi) of the theorem.

Clearly \(\mathrm{Aut}_1\mathrm{G}\) is open in Aut G with the topology \(\mathcal{T}_{\beta}\) . Let \(w\in \mathrm{Aut}G\) . Let \(\sigma\) be the automorphism \(v\mapsto wvw^{-1}\) of \(\mathrm{Aut}_1\mathrm{G}\) . It is real analytic (§ 8, no. 1, Theorem 1), \(\mathrm{L}(\sigma)\) is an \(\mathbf{R}\) -automorphism of \(\mathrm{L}(\mathrm{Aut}_1\mathrm{G})\) and

\[
\mathrm{D} \circ \mathrm{L} (\mathrm{Aut} _ {1} \mathrm{G}) \circ \mathrm{D} ^ {- 1}
\]

is an R-automorphism of a. This automorphism is also the automorphism of a derived from w by transport of structure; as w is K-analytic, we see that \(L(\sigma)\) is K-linear. Hence \(\sigma\) is K-analytic ( \(\S\) 3, no. 8, Proposition 32). By \(\S\) 1, no. 9, Proposition 18, there exists on Aut G one and only one Lie K-group structure such that \(Aut_1G\) is an open Lie subgroup of Aut G. The fact that this structure has the properties of the theorem follows from the corresponding properties for \(Aut_1G\) .

COROLLARY 1. Let G be a finite-dimensional real Lie group and \(G_{0}\) its identity component. Suppose that G is generated by \(G_{0}\) and a finite number of elements. Then Aut G has the topology \(T_{\beta}\) and is a finite-dimensional real Lie group.

COROLLARY 2. Let G be a semi-simple connected real or complex Lie group. The group Int G is the identity component of Aut G.

The mapping \(u \mapsto \mathrm{L}(u)\) is an isomorphism of Aut G onto a Lie subgroup of Aut L(G) (Theorem 1). The image of Int G under this isomorphism is Ad G. But Ad G is the identity component of Aut L(G) (§ 9, no. 8, Proposition 30 (ii)).

\subsection{THE AUTOMORPHISM GROUP OF A LIE GROUP (ULTRAMETRIC CASE)}

THEOREM 2. When K is ultrametric and locally compact and G is a compact Lie group, assertions (i), (ii), (iii), (v) and (vi) of Theorem 1 are true.

\begin{enumerate}
\def\labelenumi{(\alph{enumi})}
\item
  The uniqueness of the analytic structure considered in (i) is obvious.
\item
  Suppose that G is the Lie group defined by the normable Lie algebra L. Then G is an open and closed ball in L. Let \(w \in \operatorname{Aut} G\) . Then \(\operatorname{L}(w)\) coincides with w in a neighbourhood of 0. Let \(x \in G\) . Let p be the characteristic of the residue field. Then \(p^{n}x\) tends to 0 as n tends to \(+\infty\) . There therefore exists n such that \(w(p^{n}x) = \operatorname{L}(w)(p^{n}x)\) . Therefore
\end{enumerate}

\[
\begin{array}{r l} p ^ {n} w (x) & = w (x) ^ {p ^ {n}} = w (x ^ {p ^ {n}}) = w (p ^ {n} x) \\ & = \mathrm{L} (w) (p ^ {n} x) = p ^ {n} \mathrm{L} (w) (x) \end{array}
\]

whence \(w(x) = \mathrm{L}(w)(x)\) . Thus, \(w = \mathrm{L}(w)|\mathrm{G}\) .

Let \(\Gamma\) be the set of \(\gamma \in \operatorname{Aut} L(G)\) such that \(\gamma(G) = G\) . As \(G\) is open and compact in \(L(G)\) , \(\Gamma\) is an open subgroup of \(\operatorname{Aut} L(G)\) . By the above, Aut \(G\) is identified with \(\Gamma\) , whence there is a Lie group structure on Aut \(G\) , with which properties (i), (ii), (iii) and (vi) of Theorem 1 are obvious. Property (v) follows from Propositions 1 and 3 of no. 1.

\begin{enumerate}
\def\labelenumi{(\alph{enumi})}
\setcounter{enumi}{2}
\tightlist
\item
  We pass to the general case. By § 7, no. 1, Proposition 1, there exists an open compact subgroup \(G_0\) of \(G\) which is of the type considered in (b). Then \(G\) is generated by \(G_0\) and a finite number of elements \(x_1, x_2, \ldots, x_n\) . Let \(\text{Aut}_1G\) be the set of \(u \in \text{Aut}G\) such that \(u(G_0) = G_0\) and \(u(x_iG_0) = x_iG_0\) for \(1 \leqslant i \leqslant n\) . We define as in the proof of Theorem 1, part (c), a semi-direct product \(P\) of \(\text{Aut}G_0\) by \(G_0^n\) and an injective homomorphism \(\zeta\) of \(\text{Aut}_1G\) into \(P\) , whose image is closed in \(P\) .
\end{enumerate}

(d), (e): the argument is exactly as in parts (d), (e) of the proof of Theorem 1 with \(\mathbf{R}\) replaced by \(\mathbf{Q}_p\) and using Proposition 3 instead of Proposition 2.

Remark. If \(\mathbf{K} = \mathbf{Q}_p\) and the Lie group \(\mathbf{G}\) is generated by a compact subset (cf.~Exercise 2), assertions (i), (ii), (iii) and (vi) of Theorem 1 are still true, but not (v) (Exercise 3).

\appendixsection{Operations on linear representations}


Let \(G\) be a group, \(k\) a commutative field, \(E_1, E_2, \ldots, E_n\) vector spaces over \(k\) and \(\pi_i\) a linear representation of \(G\) on \(E_i\) ( \(1 \leqslant i \leqslant n\) ). The mapping

\[
g \mapsto \pi_ {1} (g) \otimes \dots \otimes \pi_ {n} (g)
\]

is a linear mapping of G into the vector space \(E_{1} \otimes \cdots \otimes E_{n}\) , called the tensor product of \(\pi_{1}, \ldots, \pi_{n}\) and denoted by \(\pi_{1} \otimes \cdots \otimes \pi_{n}\) .

Let E be a vector space over k and \(\pi\) a linear representation of G on E. For all \(g \in G\) , let \(\tau(g)\) (resp. \(\sigma(g)\) , \(\varepsilon(g)\) ) be the unique automorphism of the algebra \(\mathsf{T}(\mathsf{E})\) (resp. \(\mathsf{S}(\mathsf{E})\) , \(\bigwedge(\mathsf{E})\) ) which extends \(\pi(g)\) (Algebra, Chapter III, § 5, no. 2, § 6, no. 2 and § 7, no. 2). Then \(\tau\) (resp. \(\sigma\) , \(\varepsilon\) ) is a linear representation of G on T(E) (resp. \(\mathsf{S}(\mathsf{E})\) , \(\bigwedge(\mathsf{E})\) ) denoted by \(\mathsf{T}(\pi)\) (resp. \(\mathsf{S}(\pi)\) , \(\bigwedge(\pi)\) ). The subrepresentation of T(π) (resp. \(\mathsf{S}(\pi)\) , \(\bigwedge(\pi)\) ) defined by \(\mathsf{T}^{\mathsf{n}}(\mathsf{E})\) (resp. \(\mathsf{S}^{\mathsf{n}}(\mathsf{E})\) , \(\bigwedge^{\mathsf{n}}(\mathsf{E})\) ) is called the n-th tensor (resp. symmetric, exterior) power of \(\pi\) and is denoted by \(\mathsf{T}^{\mathsf{n}}(\pi)\) (resp. \(\mathsf{S}^{\mathsf{n}}(\pi)\) , \(\bigwedge^{\mathsf{n}}(\pi)\) ). Then

\[
T ^ {n} (\pi) = \pi \otimes \pi \otimes \dots \otimes \pi
\]

(n factors). The representations \(\mathsf{S}(\pi), \bigwedge (\pi)\) are quotient representations of \(\mathsf{T}(\pi)\) and hence \(\mathsf{S}^n (\pi), \bigwedge^n (\pi)\) are quotient representations of \(\mathsf{T}^n (\pi)\) .

Let g be a Lie algebra over k. The tensor product of a finite number of representations of g has already been defined in Chapter I, §3, no. 2; it is denoted by \(\pi_{1}\otimes\cdots\otimes\pi_{n}\) . Let E be a vector space over k and \(\pi\) a representation of g on E. For all \(x\in g\) , let \(\tau'(x)\) (resp. \(\sigma'(x)\) , \(\varepsilon'(x)\) ) be the unique derivation of the algebra \(\mathrm{T(E)}\) (resp. \(\mathrm{S(E)}\) , \(\bigwedge(\mathrm{E})\) ) which extends \(\pi(x)\) (Algebra, Chapter III, §10, no. 9, Example 1). Then \(\tau'\) (resp. \(\sigma'\) , \(\varepsilon'\) ) is a linear representation of g on T(E) (resp. S(E), \(\bigwedge(\mathrm{E})\) ) by Algebra, loc. cit., formula (35),

denoted by \(\mathsf{T}(\pi)\) (resp. \(\mathsf{S}(\pi),\bigwedge (\pi)\) ). The subrepresentation of \(\mathsf{T}(\pi)\) (resp. \(\mathsf{S}(\pi),\bigwedge (\pi)\) ) defined by \(\mathsf{T}^n (\mathsf{E})\) (resp. \(\mathsf{S}^n (\mathsf{E}),\bigwedge^n (\mathsf{E}))\) is denoted by \(\mathsf{T}^n (\pi)\) (resp. \(\mathsf{S}^n (\pi),\bigwedge^n (\pi))\) . The representation \(\mathsf{T}^n (\pi)\) is the tensor product on \(n\) representations identical with \(\pi\) . The representations \(\mathsf{S}(\pi),\bigwedge (\pi)\) are quotient representations of \(\mathsf{T}(\pi)\) and hence \(\mathsf{S}^n (\pi),\bigwedge^n (\pi)\) are quotient representations of \(\mathsf{T}^n (\pi)\) .

\specialchapter{Exercises}

\exercisesfor{1}

\begin{enumerate}
\def\labelenumi{\arabic{enumi}.}
\tightlist
\item
  Let G be a finite-dimensional connected real or complex Lie group. Let H and L be Lie subgroups of G such that HL = G. Show that the canonical mappings \(G \rightarrow G/H\) and \(G \rightarrow G/L\) define an isomorphism of analytic manifolds
\end{enumerate}

\[
\mathrm{G} / (\mathrm{H} \cap \mathrm{L}) \rightarrow (\mathrm{G} / \mathrm{H}) \times (\mathrm{G} / \mathrm{L}).
\]

\begin{enumerate}
\def\labelenumi{\arabic{enumi}.}
\setcounter{enumi}{1}
\item
  Let \(\mathrm{P} = \{z\in \mathbf{C}|\mathcal{I}(z) > 0\}\) . Let \(\mathbf{G}\) be the group of bijections \(z\mapsto az + b\) of \(\mathrm{P}(a > 0,b\in \mathbf{R})\) . Then \(\mathbf{G}\) operates on \(\mathrm{P}\) simply transitively, which allows us to transport the complex analytic manifold structure on \(\mathbf{P}\) to \(\mathbf{G}\) . Show that the structure obtained is invariant under left translations of \(\mathbf{G}\) but not under right translations.
\item
  Let \(R_{d}\) be the additive group of real numbers with the discrete manifold structure. Let \(R_{d}\) operate on the analytic manifold R by the law of operation \((x,y)\mapsto x+y\) . Then \(R_{d}\) operates transitively on R but R is not a Lie homogeneous space of \(R_{d}\) .
\item
  Let H be a compact real Lie group operating on a finite-dimensional real manifold V; suppose that V is of class \(C^{r}\) , where \(r \in N_{R}\) , and that the action of H on V is of class \(C^{r}\) . Let \(p \in V\) be invariant under H; the group H operates linearly on the tangent space T to V at p.~Show that there exist an open neighbourhood U of p which is stable under H and a \(C^{r}\) -morphism \(f: U \to T\) such that:
\end{enumerate}

\begin{enumerate}
\def\labelenumi{(\alph{enumi})}
\item
  \(f(p) = 0\) and the tangent mapping to \(f\) at \(p\) is the identity;
\item
  \(f\) commutes with the action of \(\mathbf{H}\) on \(\mathbf{U}\) and \(\mathbf{T}\) .
\end{enumerate}

(Choose first an \(f_0\) satisfying \((a)\) and then define \(f\) by the formula

\[
f (x) = \int_ {\mathbf {H}} h. f _ {0} (h ^ {- 1} x) d h,
\]

where \(dh\) is the Haar measure on \(H\) with total mass 1.)

Deduce that the H-spaces V and T are locally isomorphic to a neighbourhood of p, in other words that there exist systems of coordinates on V at p with respect to which H operates linearly (``Bochner's Theorem'').

\begin{enumerate}
\def\labelenumi{\arabic{enumi}.}
\setcounter{enumi}{4}
\tightlist
\item
  Perform Exercise 4 for manifolds over an ultrametric field K, assuming that H is a finite group whose order n is such that \(n.1 \neq 0\) in K.
\end{enumerate}

¶ 6. Let G be a real Lie group operating properly on a finite-dimensional Hausdorff real analytic manifold V. Let \(p \in V\) , let H be the stabilizer of \(p\) and let Gp be the orbit of \(p\) . The group H is compact and the manifold Gp is identified with the Lie homogeneous space G/H.

\begin{enumerate}
\def\labelenumi{(\alph{enumi})}
\item
  Show that there exists a submanifold S of V, passing through p, stable under H and such that \(\mathrm{T}_{p}(\mathrm{V})\) is the direct sum of \(\mathrm{T}_{p}(\mathrm{G}p)\) and \(\mathrm{T}_{p}(\mathrm{S})\) . (Choose a supplement of \(\mathrm{T}_{p}(\mathrm{G}p)\) in \(\mathrm{T}_{p}(\mathrm{V})\) which is stable under H and apply Exercise 4.)
\item
  Suppose that S is chosen as above. The group H operates properly and freely on G × S by \(h.(g,s)=(gh,h^{-1}s)\) ; let \(\mathrm{E}=(\mathrm{G}\times\mathrm{S})/\mathrm{H}\) be the corresponding quotient manifold (cf.~Differentiable and Analytic Manifolds, R, 6.5.1). The mapping \((g,s)\mapsto gs\) defines when passing to the quotient a morphism \(\rho:E\to V\) . Show that \(\rho\) commutes with the action of G on E and V and that there exists an open submanifold \(S'\) of S containing p, stable under H and such that \(\rho\) induces a manifold isomorphism of \(\mathrm{E}'=(\mathrm{G}\times\mathrm{S}')/\mathrm{H}\) onto a saturated neighbourhood of the orbit Gp.
\item
  Let \(N_{p} = T_{p}(V)/T_{p}(Gp)\) be the transversal space at p to Gp (Differentiable and Analytic Manifolds, R, 5.8.8); the group H operates on \(N_{p}\) . Deduce from (b) that knowing H and the linear representation of H on \(N_{p}\) determines the G-space V in a neighbourhood of the orbit of p.
\item
  If \(x \in \mathbf{N}_p\) , let \(\mathbf{H}_x\) be the stabilizer of \(x\) in \(\mathbf{H}\) . Show that there exists a saturated neighbourhood \(V'\) of \(Gp\) with the following property:
\end{enumerate}

For all \(p' \in V'\) , there exists \(x \in \mathbb{N}_p\) such that the stabilizer of \(p'\) in \(G\) is conjugate (in \(G\) ) to \(H_x\) .

\begin{enumerate}
\def\labelenumi{(\alph{enumi})}
\setcounter{enumi}{4}
\tightlist
\item
  Show that the \(H_{x}\) are finite in number, up to conjugation in H. (Argue by induction on dim \(N_{p}=n\) , using the action of H on a sphere of dimension n-1 which is stable under H.)†
\end{enumerate}

\begin{enumerate}
\def\labelenumi{\arabic{enumi}.}
\setcounter{enumi}{6}
\tightlist
\item
  Let E be a complete normable space over R and F a closed vector subspace of E such that there exists no bicontinuous linear bijection of E onto \(\mathbf{F} \times (\mathbf{E}/\mathbf{F})\) . Let X be the analytic manifold \(R \times E \times F\) . Let G be the additive group of F, which operates on X by the mapping
\end{enumerate}

\[
(g, (\lambda , e, f)) \mapsto (\lambda , e + g, f + \lambda g)
\]

of \(\mathbf{G} \times \mathbf{X}\) into \(\mathbf{X}\) . For all \(x \in \mathbf{X}\) , let \(\mathbf{R}_x\) be the G-orbit of \(x\) , which is a quasi-submanifold of \(\mathbf{X}\) . Then condition (a) of Proposition 10 is satisfied, but there exists no ordered pair \((\mathbf{Y}, \pi)\) with the following properties: \(\mathbf{Y}\) is an analytic manifold, \(\pi\) is a morphism of \(\mathbf{X}\) into \(\mathbf{Y}\) and, for all \(x \in \mathbf{X}\) ,

\[
\mathrm{T} _ {x} (\pi): \mathrm{T} _ {x} (\mathrm{X}) \rightarrow \mathrm{T} _ {\pi (x)} (\mathrm{Y})
\]

is surjective with kernel \(\mathrm{T}_x(\mathbf{R}_x)\) .

(Suppose that \((\mathrm{Y},\pi)\) exists. Let \(\mathbf{H}\) be the tangent space to \(\mathbf{Y}\) at \(\pi (0,0,0)\) . Then \(\mathbf{H}\) is isomorphic to \(\mathbf{R}\times \mathbf{F}\times (\mathrm{E} / \mathrm{F})\) . On the other hand \(\mathbf{H}\) is isomorphic to \(\mathrm{T}_{\pi (x)}\mathrm{Y}\) for \(x\) sufficiently close to \((0,0,0)\) and is hence isomorphic to \(\mathbf{R}\times \mathbf{E}\) .)

\begin{enumerate}
\def\labelenumi{\arabic{enumi}.}
\setcounter{enumi}{7}
\item
  Let \(\mathbf{T}\) be given the normalized Haar and \(\mathcal{L}(\mathrm{L}^2 (\mathbf{T}))\) the normic topology. Show that the regular representation of \(\mathbf{T}\) on \(\mathrm{L}^2 (\mathbf{T})\) is not continuous. (For all \(g\in \mathbf{T}\) distinct from \(\epsilon\) , construct a function \(f\in \mathrm{L}^2 (\mathbf{T})\) such that \(\| f\| = 1\) , \(\| \gamma (g)f - f\| = \sqrt{2}\) .)
\item
  Let I be the set of \((x, \sqrt{2} x) \in \mathbf{R}^2\) for \(-\frac{1}{2} < x < \frac{1}{2}\) . Let U be the canonical image of I in the real Lie group \(\mathbf{H} = \mathbf{T}^2\) . Show that U is a Lie subgroup germ of H but that the subgroup \(\mathrm{H}'\) of H generated by U is dense in H and not closed in H.
\item
  Let \(\mathbf{R}_d\) be the group \(\mathbf{R}\) with the discrete manifold structure. Let \(\mathrm{H}\) be the real Lie group \(\mathbf{R} \times \mathbf{R}_d\) . Let \(\mathrm{U}\) be the set of elements of \(\mathrm{H}\) of the form \((x, x)\) or of the form \((x, -x)\) where \(x \in \mathbf{R}\) . Then \(\mathrm{U}\) is a Lie subgroup germ of \(\mathrm{H}\) , \(\mathrm{U}\) is discrete, the subgroup \(\mathrm{G}\) of \(\mathrm{H}\) generated by \(\mathrm{U}\) is equal to \(\mathrm{H}\) and the Lie group structure on \(\mathrm{G}\) defined by the Corollary to Proposition 22 is the discrete group structure.
\end{enumerate}

\exercisesfor{3}

\begin{enumerate}
\def\labelenumi{\arabic{enumi}.}
\item
  Let \(G\) be a Lie group and \(G_1\) and \(G_2\) Lie subgroups of \(G\) . Suppose that \(G_2\) is of finite codimension in \(G\) and that \(K\) is of characteristic 0. Then \(G_1 \cap G_2\) is a Lie subgroup of \(G\) with Lie algebra \(L(G_1) \cap L(G_2)\) (argue as in Proposition 29 and its Corollary 2).
\item
  Suppose that \(\mathbf{K}\) is locally compact. Let \(\mathbf{G}\) be a Lie group of finite dimension \(n\) .
\end{enumerate}

\begin{enumerate}
\def\labelenumi{(\alph{enumi})}
\tightlist
\item
  Let \(\phi\) be an étale endomorphism of G with finite kernel. Let \(\mu\) be a left Haar measure of G. Then \(\phi\) is proper and
\end{enumerate}

\[
\phi (\mu) = \alpha \cdot \mu | \phi (G) \quad \text { where } \quad \alpha = \frac {\operatorname{Card} (\operatorname{Ker} \phi)}{\operatorname{mod} \det L (\phi)}.
\]

(Use Proposition 55.)
